\section{Phase patterns and angular rules}\label{pc:sec}
\begingroup\sloppy

Fix a centre $c\in\R^2$ and three directions $\alpha<\beta<\delta<\alpha+2\pi$, read counterclockwise around
$c$. Together with the three antipodal directions $\alpha+\pi$, $\beta+\pi$, $\delta+\pi$ they give six points on
the circle of directions, distinct unless one of the gaps $\beta-\alpha$, $\delta-\beta$, $\alpha+2\pi-\delta$
equals $\pi$, in which case a direction and an antipode coincide (Remark~\ref{pc:rem:reading}). This section gives
written proofs of the Lean-checked results that classify the cyclic pattern of these points, expressed by the arc
predicates of Definition~\ref{pc:def:arc} (which are defined in every case), and connect it with the sign of
$\orient$ for three points in polar form about $c$.
\begin{itemize}
\item The pattern alternates (every open arc between consecutive directions contains an antipode; pattern
  $PAPAPA$) if and only if all three gaps are smaller than $\pi$ (Theorem~\ref{pc:thm:alternating}).
\item A run $PPP$ with middle point $\beta$ (no antipode in the arcs $(\alpha,\beta)$ and $(\beta,\delta)$) is
  equivalent to $\delta-\alpha<\pi$, provided $\alpha+2\pi-\delta\neq\pi$ (Theorem~\ref{pc:thm:runB}).
\item If no gap equals $\pi$, the pattern alternates or has a run, and alternation excludes every run
  (Theorem~\ref{pc:thm:trichotomy}).
\item Rule~A: an alternating triple at positive radii turns counterclockwise (Theorem~\ref{pc:thm:rule-alt}).
  Rule~B: a triple with a run turns counterclockwise if the gap not adjacent to the middle point of the run
  differs from $\pi$, and the other two radii are positive and at most the radius of the middle point
  (Theorem~\ref{pc:thm:rule-run}). The two rules combined (Theorem~\ref{pc:thm:phase}), and a version for
  arbitrary real directions (Theorem~\ref{pc:thm:sign}).
\item The dichotomy for cyclic Boolean words of length $6$ (Theorem~\ref{pc:thm:word}; computer-assisted in
  Lean, with a written proof by hand here).
\item Rule~B on the convex layers of a rotationally symmetric set, where the radius condition follows from the
  layer indices (Theorem~\ref{pc:thm:layers}).
\end{itemize}

\paragraph{Sources.} File paths are relative to the results directory \dirn{cloud/es7-2026-10-19/results} of
the repository. The Lean sources are:
\begin{itemize}
\item \dirn{LEANEXT/phase/PhaseClassification.lean} (545 lines, SHA-256 \texttt{490efcd0a7e4}\dots), namespace
  \lean{PhaseClassification} (\dirn{LEANEXT/phase/REPORT.md} gives 540 lines for this file, with the same hash
  prefix);
\item \dirn{LEANEXT/linked/Linked.lean} (1416 lines, SHA-256 \texttt{f1e18c119276}\dots); its new lines are
  1385--1416 (\lean{dist2_polar} and \lean{orient_pos_of_layers}, namespace \lean{ErdosSzekeres.LayerLemma});
\item \dirn{LEANEXT/capstone/Capstone.lean} (2004 lines, SHA-256 \texttt{dc382395ca55}\dots); its new lines are
  1966--2004 (namespace \lean{ErdosSzekeres.Capstone});
\item \dirn{LEANEXT/capstone/Capstone2.lean} (2132 lines, SHA-256 \texttt{f8c3930c2427}\dots); its new lines are
  2004--2132.
\end{itemize}
The following relations were checked with \texttt{diff}. Lines 40--102 of \dirn{PhaseClassification.lean}
repeat lines 931--996 of \dirn{LAYERLEAN/LayerLemma.lean} (\lean{polar}, \lean{orient_polar},
\lean{orient_polar_centre_inside}, \lean{sin_add_lt}, \lean{orient_polar_halfplane}) as new constants of the
namespace \lean{PhaseClassification}; only the docstrings differ. Lines 1--1147 of \dirn{Linked.lean} are
lines 1--1147 of \dirn{LAYERLEAN/LayerLemma.lean}, and lines 1--1384 of \dirn{Linked.lean} are those of
\dirn{LEANEXT/merged/Merged.lean}. \dirn{Capstone.lean} consists of an 18-line header, lines 11--1416 of
\dirn{Linked.lean} (as its lines 19--1424), a blank line, lines 6--545 of \dirn{PhaseClassification.lean} (as
its lines 1426--1965) and the new lines 1966--2004. Lines 1--2003 of \dirn{Capstone2.lean} are those of
\dirn{Capstone.lean}. Hence a declaration at line $L$ of \dirn{PhaseClassification.lean} reappears at line
$L+1420$ of \dirn{Capstone.lean} and \dirn{Capstone2.lean}, and one at line $L\ge 11$ of \dirn{Linked.lean}
reappears at line $L+8$. We cite every declaration at its first occurrence. \dirn{Capstone2.lean} is
byte-identical to the copy \dirn{ENCBLEAN/geo/ES8Geo/Capstone2.lean} that the ENCBLEAN geometric package
imports (hash list \dirn{ENCBLEAN/geo/SOURCES.sha256}). Further Lean files cited below, with the first 12 hex
digits of their SHA-256: \dirn{LAYERLEAN/LayerLemma.lean} (\texttt{afb8542dceb0}), \dirn{LEANEXT/merged/Merged.lean}
(\texttt{245b7d65407a}), \dirn{ENCBLEAN/geo/ES8Geo/Audit.lean} (\texttt{39f214bbf30b}),
\dirn{ENCBLEAN/geo/ES8Geo/Labelling.lean} (\texttt{489ba89f8d6f}), \dirn{ENCBLEAN/geo/ES8Geo/ArcBridge.lean}
(\texttt{46966b48a101}), \dirn{ENCBLEAN/geo/ES8Geo/ModelOfLabelling.lean} (\texttt{5b532f49cc11}) and
\dirn{ENCBLEAN/EncB/SoundWord.lean} (\texttt{c1bdb4175ee5}). The files were built with Lean 4.24.0 and, except
\dirn{ENCBLEAN/EncB/SoundWord.lean} (part of a package without Mathlib), with Mathlib \texttt{f897ebcf}; these
are the repository's pins.

\paragraph{Rigor and axioms.} Every definition, lemma, proposition, theorem and corollary below whose head or one
of whose parts carries a Lean reference is \emph{Lean-checked}; the proofs given here are written proofs (not
refereed). Remarks are commentary and, even when they name Lean declarations, are not Lean-checked as a whole; the
mathematical claims they add to the Lean statements are written steps and are not Lean-checked. The log
\dirn{LEANEXT/phase/logs/axioms.log} records the output of \texttt{\#print axioms} for 16 theorems of
\dirn{PhaseClassification.lean} (\lean{alternating_iff_gaps_lt_pi}, \lean{alternating_exactly_one},
\lean{runMidB_iff}, \lean{gap_classification}, \lean{phase_trichotomy}, \lean{orient_pos_of_alternating},
\lean{orient_pos_of_runMidA}, \lean{orient_pos_of_runMidB}, \lean{orient_pos_of_runMidD},
\lean{orient_pos_of_phase}, \lean{orient_sign_of_phase}, \lean{centre_inside_of_alternating},
\lean{halfplane_of_runMidB}, \lean{exists_polar}, \lean{polar_toIcoMod}, \lean{word_dichotomy}); each depends
only on \lean{propext}, \lean{Classical.choice} and \lean{Quot.sound}. Every other declaration of that file
covered here is used in the statement or the proof of one of these 16 theorems, so its axioms are among these
three. For the new
declarations of \dirn{Linked.lean}, \dirn{Capstone.lean} and \dirn{Capstone2.lean} no axiom output is saved in
\dirn{LEANEXT}; \dirn{LEANEXT/linked/REPORT.md} and \dirn{LEANEXT/capstone/REPORT.md} state the same three
axioms. Of these declarations, \lean{dist2_polar}, \lean{Capstone.polar_eq}, \lean{Capstone.exists_angle},
\lean{Capstone.sqrt_dist2_pos} and \lean{Capstone.ne_of_mem} (with \lean{Capstone.polar_zero}, used in the
proof of \lean{sqrt_dist2_pos}) are used, together with \lean{PhaseClassification.orient_pos_of_alternating},
\lean{PhaseClassification.orient_pos_of_runMidB} and \lean{PhaseClassification.polar_add_int_two_pi}, through the
byte-identical copy of \dirn{Capstone2.lean}, in proofs in \dirn{ENCBLEAN/geo/ES8Geo/Labelling.lean} and
\dirn{ENCBLEAN/geo/ES8Geo/ArcBridge.lean} on which the declarations \lean{EncB.Geo.labelling_of_counterexample}
(in \dirn{Labelling.lean}) and \lean{EncB.Geo.modelOfLabelling} (in
\dirn{ENCBLEAN/geo/ES8Geo/ModelOfLabelling.lean}) depend. The axioms of these two declarations are pinned to the
same three by \texttt{\#guard\_msgs} in \dirn{ENCBLEAN/geo/ES8Geo/Audit.lean}, lines 45--49. The package compiles
the copy from the same sources, against the same imports, as the original file: besides the byte-identical copy
and the common Lean and Mathlib pins (see Sources), the ten repository modules that \dirn{Capstone2.lean} needs
(its imports \lean{ErdosSzekeres.Definitions}, \lean{ErdosSzekeres.Orientation},
\lean{ErdosSzekeres.OrientationGeometry} and \lean{ErdosSzekeres.Suk.OrientationClique}, with the repository
modules they import) are compiled there from verbatim copies, whose SHA-256 hashes are listed in
\dirn{ENCBLEAN/geo/VENDOR.sha256} and agree with those of the repository files. So the axioms of the declarations
used are among these three as well. The remaining new declarations (\lean{orient_pos_of_layers},
the six layer rules of \lean{ErdosSzekeres.Capstone} and \lean{Capstone.orient_pos_of_alternating_polar}) are not
used in the ENCBLEAN geometric package outside the copy itself, and no \texttt{\#print axioms} output is saved
for them; besides the axiom statements of the two
reports, the evidence for them is the clean builds stated in these reports and the ENCBLEAN build log
\dirn{ENCBLEAN/logs/m4_build.log}, a successful build of the package that imports the byte-identical copy of
\dirn{Capstone2.lean} as a module, in which only the \lean{KeyOrder} module prints messages (its axiom lines). A
text search of the four files finds no \lean{sorry}, \lean{admit}, \lean{axiom}
declaration or \lean{native_decide}; among the declarations of this section, \lean{decide} occurs only in the
proof of \lean{word_dichotomy}.

\paragraph{Dependencies.} We use Definitions~\ref{fd:def:point}, \ref{fd:def:orient}, \ref{fd:def:rot},
\ref{fd:def:inv} and \ref{fd:def:ngon}, and Lemma~\ref{fd:lem:orient}. From the layer sections we use the hull
vertices $\hv$ and the peeling $\peel$ (Definitions~\ref{ll:def:hv} and~\ref{ll:def:peel}), the constant
\lean{LayerLemma.polar} (Definition~\ref{lg:def:polar}), its polar formula and angular rules
(Lemmas~\ref{lg:lem:polar}, \ref{lg:lem:ruleA}, \ref{lg:lem:sinadd} and~\ref{lg:lem:ruleB}), and the radius
order of the layers (Theorem~\ref{lg:thm:radius}).

\subsection{Conventions and polar points}

Points are elements of $\R^2$, written $p=(p_0,p_1)$ (Lean \lean{Point}, Definition~\ref{fd:def:point}); two
points are equal when both coordinates agree. We use
$\orient(a,b,d)=(b_0-a_0)(d_1-a_1)-(b_1-a_1)(d_0-a_0)$ (Definition~\ref{fd:def:orient}), which is positive when
$a,b,d$ turn counterclockwise, and from Lemma~\ref{fd:lem:orient}: $\orient(a,d,b)=-\orient(a,b,d)$
(\lean{orient_swap}), $\orient(b,d,a)=\orient(a,b,d)$ (\lean{orient_cycle}) and
$\orient(a+v,b+v,d+v)=\orient(a,b,d)$ (\lean{orient_translate}). Angles are arbitrary real numbers. The number
$\pi>0$ is Lean's \lean{Real.pi}, and an integer $k$ is identified with its image in $\R$ (Lean writes this cast
explicitly).

\emph{Reduction modulo $2\pi$.} For $\alpha,\theta\in\R$ we write $\operatorname{toIcoMod}(\alpha,\theta)$ for
Mathlib's \lean{toIcoMod} with period $2\pi$ (its positivity argument is \lean{Real.two_pi_pos}), base point
$\alpha$ and argument $\theta$. By definition $\operatorname{toIcoMod}(\alpha,\theta)=\theta-2\pi n$, where
$n=\operatorname{toIcoDiv}(\alpha,\theta)$ is the unique integer with $\theta-2\pi n\in[\alpha,\alpha+2\pi)$,
namely $n=\lfloor(\theta-\alpha)/(2\pi)\rfloor$. Only two properties are used (Mathlib \lean{toIcoMod_mem_Ico}
and \lean{toIcoMod_add_toIcoDiv_zsmul}):
\begin{enumerate}[label=(T\arabic*)]
\item $\alpha\le\operatorname{toIcoMod}(\alpha,\theta)<\alpha+2\pi$;
\item $\theta=\operatorname{toIcoMod}(\alpha,\theta)+n\cdot 2\pi$ with the integer
  $n=\operatorname{toIcoDiv}(\alpha,\theta)$.
\end{enumerate}

\begin{definition}[polar points and squared distance; \leanref{PhaseClassification.polar}{LEANEXT/phase/PhaseClassification.lean:41}, \leanref{ErdosSzekeres.LayerLemma.dist2}{LAYERLEAN/LayerLemma.lean:999}]\label{pc:def:polar}
For $c\in\R^2$ and $r,\theta\in\R$,
\[
  \operatorname{polar}(c,r,\theta)=(c_0+r\cos\theta,\ c_1+r\sin\theta).
\]
No sign condition is imposed on $r$. For $c,y\in\R^2$, $\operatorname{dist2}(c,y)=(y_0-c_0)^2+(y_1-c_1)^2$.
\end{definition}

The constant \lean{ErdosSzekeres.LayerLemma.polar} (Definition~\ref{lg:def:polar};
\dirn{LAYERLEAN/LayerLemma.lean:932}, repeated at \dirn{LEANEXT/linked/Linked.lean:932} and
\dirn{LEANEXT/capstone/Capstone2.lean:940}) is a different Lean constant given by the same formula;
Lemma~\ref{pc:lem:helpers}(a) records that the two agree. We write $\operatorname{polar}$ for both; every
statement below that involves $\operatorname{polar}$ says, directly or through the convention of its subsection,
which of the two constants its Lean version uses. The definition of $\operatorname{dist2}$ is repeated at
\dirn{LEANEXT/linked/Linked.lean:999} and \dirn{LEANEXT/capstone/Capstone2.lean:1007}.

\subsection{The polar orientation formula and the two rules}

\begin{lemma}[polar formula, rule A, rule B; \leanref{PhaseClassification.orient_polar}{LEANEXT/phase/PhaseClassification.lean:44} for (a), and the declarations cited in (b)--(e)]\label{pc:lem:polar-facts}
Let $c\in\R^2$; all other letters denote real numbers, and $\operatorname{polar}$ is
\lean{PhaseClassification.polar}.
\begin{enumerate}[label=(\alph*)]
\item $\orient(\operatorname{polar}(c,r_a,\alpha),\operatorname{polar}(c,r_b,\beta),\operatorname{polar}(c,r_d,\delta))
  =r_ar_b\sin(\beta-\alpha)+r_br_d\sin(\delta-\beta)+r_dr_a\sin(\alpha-\delta)$.
\item \textup{(\leanref{PhaseClassification.orient_polar_centre}{LEANEXT/phase/PhaseClassification.lean:407}.)}
  $\orient(\operatorname{polar}(c,r_a,\alpha),\operatorname{polar}(c,r_b,\beta),c)=r_ar_b\sin(\beta-\alpha)$.
\item \textup{(Rule A; \leanref{PhaseClassification.orient_polar_centre_inside}{LEANEXT/phase/PhaseClassification.lean:51}.)}
  If $0<r_a$, $0<r_b$, $0<r_d$, $0<\beta-\alpha$, $\beta-\alpha<\pi$, $0<\delta-\beta$,
  $\delta-\beta<\pi$, $0<\alpha-\delta+2\pi$ and $\alpha-\delta+2\pi<\pi$, then
  $0<\orient(\operatorname{polar}(c,r_a,\alpha),\operatorname{polar}(c,r_b,\beta),\operatorname{polar}(c,r_d,\delta))$.
\item \textup{(\leanref{PhaseClassification.sin_add_lt}{LEANEXT/phase/PhaseClassification.lean:64}.)}
  If $0<x$, $x<\pi$, $0<y$ and $y<\pi$, then $\sin(x+y)<\sin x+\sin y$.
\item \textup{(Rule B; \leanref{PhaseClassification.orient_polar_halfplane}{LEANEXT/phase/PhaseClassification.lean:80}.)}
  If $0<r_a$, $0<r_d$, $r_a\le r_b$, $r_d\le r_b$, $0<\beta-\alpha$, $0<\delta-\beta$ and
  $\delta-\alpha<\pi$, then
  $0<\orient(\operatorname{polar}(c,r_a,\alpha),\operatorname{polar}(c,r_b,\beta),\operatorname{polar}(c,r_d,\delta))$.
\end{enumerate}
\end{lemma}

Parts (a), (c) and (e) are the statements of Lemmas~\ref{lg:lem:polar}, \ref{lg:lem:ruleA}
and~\ref{lg:lem:ruleB} with the constant \lean{PhaseClassification.polar} in place of \lean{LayerLemma.polar}, and
part (d) is the statement of Lemma~\ref{lg:lem:sinadd}; the Lean proofs are the same. For completeness we repeat
the short arguments.

\begin{proof}
(a) Put $u=(r_a\cos\alpha,r_a\sin\alpha)$, $v=(r_b\cos\beta,r_b\sin\beta)$, $w=(r_d\cos\delta,r_d\sin\delta)$,
so that the three points are $c+u$, $c+v$, $c+w$, and write $[x,y]=x_0y_1-x_1y_0$ for $x,y\in\R^2$. By
translation invariance, $\orient(c+u,c+v,c+w)=\orient(u,v,w)=[v-u,w-u]$, and by bilinearity and antisymmetry of
$[\cdot,\cdot]$,
\[
  [v-u,w-u]=[v,w]-[v,u]-[u,w]+[u,u]=[u,v]+[v,w]+[w,u].
\]
Finally $[u,v]=r_ar_b(\cos\alpha\sin\beta-\sin\alpha\cos\beta)=r_ar_b\sin(\beta-\alpha)$, and in the same way
$[v,w]=r_br_d\sin(\delta-\beta)$ and $[w,u]=r_dr_a\sin(\alpha-\delta)$.

(b) For every $\delta$, $\operatorname{polar}(c,0,\delta)=c$. Hence (b) is (a) with $r_d=0$. (Lean expands the
definitions directly; the computation is the same.)

(c) By (a) the orientation is $r_ar_b\sin(\beta-\alpha)+r_br_d\sin(\delta-\beta)+r_dr_a\sin(\alpha-\delta)$. The
sine is positive on $(0,\pi)$, so $\sin(\beta-\alpha)>0$, $\sin(\delta-\beta)>0$ and
$\sin(\alpha-\delta)=\sin(\alpha-\delta+2\pi)>0$. All radii are positive, so all three terms are positive.

(d) By the addition formula, $\sin x+\sin y-\sin(x+y)=\sin x\,(1-\cos y)+\sin y\,(1-\cos x)$. Here $\sin x>0$
and $\sin y>0$. Moreover $\cos x<1$ and $\cos y<1$, because $\cos$ is strictly decreasing on $[0,\pi]$ with
$\cos 0=1$. So both products are positive.

(e) Put $x=\beta-\alpha>0$ and $y=\delta-\beta>0$. Then $x+y=\delta-\alpha<\pi$, so $x,y\in(0,\pi)$, and
$\alpha-\delta=-(x+y)$. By (a), the orientation $O$ in (e) satisfies
\begin{align*}
  O&=r_ar_b\sin x+r_br_d\sin y-r_dr_a\sin(x+y)\\
   &=r_a(r_b-r_d)\sin x+r_d(r_b-r_a)\sin y+r_ar_d\bigl(\sin x+\sin y-\sin(x+y)\bigr).
\end{align*}
The first two terms are $\ge0$ (as $r_a,r_d>0$, $r_b-r_d\ge0$, $r_b-r_a\ge0$, $\sin x,\sin y>0$), and the last
is $>0$ by (d) and $r_ar_d>0$. (Positivity of $r_b$ is not assumed; it follows from $r_b\ge r_a>0$.)
\end{proof}

\subsection{Periodicity and polar forms}

\begin{lemma}[periodicity; \leanref{PhaseClassification.polar_add_int_two_pi}{LEANEXT/phase/PhaseClassification.lean:106}, \leanref{PhaseClassification.polar_add_two_pi}{LEANEXT/phase/PhaseClassification.lean:111}, \leanref{PhaseClassification.polar_toIcoMod}{LEANEXT/phase/PhaseClassification.lean:485}]\label{pc:lem:periodic}
Let $c\in\R^2$ and $r,\theta,\alpha\in\R$, and let $\operatorname{polar}$ be \lean{PhaseClassification.polar}.
\begin{enumerate}[label=(\alph*)]
\item For every $k\in\Z$, $\operatorname{polar}(c,r,\theta+k\cdot2\pi)=\operatorname{polar}(c,r,\theta)$.
\item $\operatorname{polar}(c,r,\theta+2\pi)=\operatorname{polar}(c,r,\theta)$.
\item $\operatorname{polar}(c,r,\operatorname{toIcoMod}(\alpha,\theta))=\operatorname{polar}(c,r,\theta)$.
\end{enumerate}
\end{lemma}

\begin{proof}
(a) $\cos(\theta+2\pi k)=\cos\theta$ and $\sin(\theta+2\pi k)=\sin\theta$ for $k\in\Z$. (b) is (a) with $k=1$.
(c) By (T2), $\theta=\operatorname{toIcoMod}(\alpha,\theta)+n\cdot2\pi$ with $n\in\Z$, so by (a)
$\operatorname{polar}(c,r,\theta)=\operatorname{polar}(c,r,\operatorname{toIcoMod}(\alpha,\theta)+n\cdot2\pi)
=\operatorname{polar}(c,r,\operatorname{toIcoMod}(\alpha,\theta))$.
\end{proof}

\begin{lemma}[every point other than the centre has a polar form; \leanref{PhaseClassification.exists_polar}{LEANEXT/phase/PhaseClassification.lean:459}]\label{pc:lem:exists-polar}
For all $c,p\in\R^2$ with $p\neq c$ there are $r,\theta\in\R$ with $0<r$ and $p=\operatorname{polar}(c,r,\theta)$,
where $\operatorname{polar}$ is \lean{PhaseClassification.polar}.
\end{lemma}

\begin{proof}
Let $u=p-c\neq0$ and $r=\sqrt{u_0^2+u_1^2}>0$, and put $t=u_0/r\in[-1,1]$; note $1-t^2=u_1^2/r^2$. Recall that
$\arccos t\in[0,\pi]$, $\cos(\arccos t)=t$ and $\sin(\arccos t)=\sqrt{1-t^2}=\abs{u_1}/r$. If $u_1\ge0$, let
$\theta=\arccos t$; then $\cos\theta=u_0/r$ and $\sin\theta=\abs{u_1}/r=u_1/r$. If $u_1<0$, let
$\theta=-\arccos t$; then $\cos\theta=t=u_0/r$ and $\sin\theta=-\abs{u_1}/r=u_1/r$. Hence
$p=c+u=(c_0+r\cos\theta,c_1+r\sin\theta)=\operatorname{polar}(c,r,\theta)$. (The Lean proof takes
$r=\abs{z}$ and $\theta=\arg z$ for the complex number $z=u_0+iu_1$, using $\abs{z}\cos(\arg z)=\operatorname{Re}z$
and $\abs{z}\sin(\arg z)=\operatorname{Im}z$; this is the same polar representation.)
\end{proof}

\begin{lemma}[helpers for layered points; \leanref{ErdosSzekeres.Capstone.polar_eq}{LEANEXT/capstone/Capstone.lean:2001}, \leanref{ErdosSzekeres.Capstone.polar_zero}{LEANEXT/capstone/Capstone2.lean:2051}, \leanref{ErdosSzekeres.LayerLemma.dist2_polar}{LEANEXT/linked/Linked.lean:1390}, \leanref{ErdosSzekeres.Capstone.exists_angle}{LEANEXT/capstone/Capstone2.lean:2056}, \leanref{ErdosSzekeres.Capstone.sqrt_dist2_pos}{LEANEXT/capstone/Capstone2.lean:2064}, \leanref{ErdosSzekeres.Capstone.ne_of_mem}{LEANEXT/capstone/Capstone2.lean:2073}]\label{pc:lem:helpers}
Let $c,p\in\R^2$ and $r,\theta\in\R$. In \textup{(b)--(e)}, $\operatorname{polar}$ is the constant
\lean{LayerLemma.polar}.
\begin{enumerate}[label=(\alph*)]
\item $\text{\lean{LayerLemma.polar}}(c,r,\theta)=\text{\lean{PhaseClassification.polar}}(c,r,\theta)$.
\item $\operatorname{polar}(c,0,\theta)=c$.
\item $\operatorname{dist2}(c,\operatorname{polar}(c,r,\theta))=r^2$.
\item If $p\neq c$, there is $\theta\in\R$ with $p=\operatorname{polar}\bigl(c,\sqrt{\operatorname{dist2}(c,p)},\theta\bigr)$.
\item If $p\neq c$, then $0<\sqrt{\operatorname{dist2}(c,p)}$.
\item If $S$ is a finite set of points with $c\notin S$ and $p\in S$, then $p\neq c$.
\end{enumerate}
\end{lemma}

\begin{proof}
(a) Both constants are defined by the formula of Definition~\ref{pc:def:polar} (in Lean the proof is
\lean{rfl}). (b) $c_0+0\cdot\cos\theta=c_0$ and $c_1+0\cdot\sin\theta=c_1$.
(c) $(r\cos\theta)^2+(r\sin\theta)^2=r^2(\cos^2\theta+\sin^2\theta)=r^2$.
(d) By Lemma~\ref{pc:lem:exists-polar} and (a) there are $r>0$ and $\theta$ with $p=\operatorname{polar}(c,r,\theta)$.
By (c), $\operatorname{dist2}(c,p)=r^2$, and $\sqrt{r^2}=r$ because $r\ge0$. Hence
$p=\operatorname{polar}(c,\sqrt{\operatorname{dist2}(c,p)},\theta)$.
(e) The square root is $\ge0$. If it were $0$, then by (d) and (b) $p=\operatorname{polar}(c,0,\theta)=c$, a
contradiction.
(f) If $p=c$, then $c=p\in S$, contrary to $c\notin S$.
\end{proof}

\subsection{Arcs and phase patterns}

\begin{definition}[arcs and phase patterns; \leanref{PhaseClassification.inArc}{LEANEXT/phase/PhaseClassification.lean:117}, \leanref{PhaseClassification.arcHasAntipode}{LEANEXT/phase/PhaseClassification.lean:120}, \leanref{PhaseClassification.Alternating}{LEANEXT/phase/PhaseClassification.lean:126}, \leanref{PhaseClassification.RunMidB}{LEANEXT/phase/PhaseClassification.lean:132}, \leanref{PhaseClassification.RunMidD}{LEANEXT/phase/PhaseClassification.lean:136}, \leanref{PhaseClassification.RunMidA}{LEANEXT/phase/PhaseClassification.lean:140}]\label{pc:def:arc}
For $s,g,\theta\in\R$,
\[
  \operatorname{inArc}(s,g,\theta)\ :\Longleftrightarrow\ \exists k\in\Z:\ s<\theta+2\pi k\ \text{ and }\ \theta+2\pi k<s+g .
\]
For $\alpha,\beta,\delta,s,g\in\R$,
\[
  \operatorname{arcHasAntipode}(\alpha,\beta,\delta,s,g)\ :\Longleftrightarrow\
  \operatorname{inArc}(s,g,\alpha+\pi)\vee\operatorname{inArc}(s,g,\beta+\pi)\vee\operatorname{inArc}(s,g,\delta+\pi).
\]
For $\alpha,\beta,\delta\in\R$ abbreviate (this notation is not a Lean name)
\begin{align*}
  H_1&=\operatorname{arcHasAntipode}(\alpha,\beta,\delta,\alpha,\beta-\alpha),\\
  H_2&=\operatorname{arcHasAntipode}(\alpha,\beta,\delta,\beta,\delta-\beta),\\
  H_3&=\operatorname{arcHasAntipode}(\alpha,\beta,\delta,\delta,\alpha+2\pi-\delta),
\end{align*}
and write $H_j(\alpha,\beta,\delta)$ when the triple has to be shown. Then
\begin{align*}
  \operatorname{Alternating}(\alpha,\beta,\delta)&:\Longleftrightarrow H_1\wedge H_2\wedge H_3,&
  \operatorname{RunMidB}(\alpha,\beta,\delta)&:\Longleftrightarrow \neg H_1\wedge\neg H_2,\\
  \operatorname{RunMidD}(\alpha,\beta,\delta)&:\Longleftrightarrow \neg H_2\wedge\neg H_3,&
  \operatorname{RunMidA}(\alpha,\beta,\delta)&:\Longleftrightarrow \neg H_3\wedge\neg H_1 .
\end{align*}
\end{definition}

\begin{remark}[reading the definitions]\label{pc:rem:reading}
For $0<g\le2\pi$, $\operatorname{inArc}(s,g,\theta)$ says that the direction $\theta$, taken modulo $2\pi$, lies
on the open counterclockwise arc from $s$ to $s+g$. When $\alpha<\beta<\delta<\alpha+2\pi$, the arcs of
$H_1,H_2,H_3$ are the open arcs $(\alpha,\beta)$, $(\beta,\delta)$, $(\delta,\alpha+2\pi)$ between
counterclockwise-consecutive directions (for $H_3$, $s+g=\delta+(\alpha+2\pi-\delta)=\alpha+2\pi$), and $H_j$ says
that the $j$-th arc contains the direction of an antipode. We use the gaps $g_1=\beta-\alpha$,
$g_2=\delta-\beta$, $g_3=\alpha+2\pi-\delta$; they satisfy $g_1+g_2+g_3=2\pi$, and $g_1,g_2,g_3>0$ when
$\alpha<\beta<\delta<\alpha+2\pi$. Note $g_1+g_2=\delta-\alpha$, $g_2+g_3=\alpha+2\pi-\beta$ and
$g_3+g_1=\beta+2\pi-\delta$. If $\alpha<\beta<\delta<\alpha+2\pi$ and no gap equals $\pi$, the six directions
$\alpha,\beta,\delta,\alpha+\pi,\beta+\pi,\delta+\pi$ are distinct modulo $2\pi$ (a direction and an antipode
coincide modulo $2\pi$ exactly when $\beta-\alpha$, $\delta-\beta$ or $\delta-\alpha$ equals $\pi$, that is, when
$g_1$, $g_2$ or $g_3$ equals $\pi$). In that case, writing $P$ for a direction and $A$ for an antipode,
$\operatorname{Alternating}$ is the cyclic pattern $PAPAPA$, and $\operatorname{RunMidB}$ (resp.\
$\operatorname{RunMidD}$, $\operatorname{RunMidA}$) is a cyclic run $PPP$ whose middle point is $\beta$ (resp.\
$\delta$, $\alpha$). The definitions themselves assume no order; every result below states the order it uses.
\end{remark}

\begin{lemma}[membership and non-membership; \leanref{PhaseClassification.inArc_of}{LEANEXT/phase/PhaseClassification.lean:143}, \leanref{PhaseClassification.not_inArc}{LEANEXT/phase/PhaseClassification.lean:147}]\label{pc:lem:inarc}
Let $s,g,\theta\in\R$.
\begin{enumerate}[label=(\alph*)]
\item For every $k\in\Z$: if $s<\theta+2\pi k$ and $\theta+2\pi k<s+g$, then $\operatorname{inArc}(s,g,\theta)$.
\item For every $m\in\Z$: if $s+g\le\theta+2\pi m$ and $\theta+2\pi m\le s+2\pi$, then
  $\neg\operatorname{inArc}(s,g,\theta)$.
\end{enumerate}
\end{lemma}

\begin{proof}
(a) is the definition, with witness $k$.
(b) Suppose $k\in\Z$ satisfies $s<\theta+2\pi k<s+g$. If $k\ge m$, then $\theta+2\pi k\ge\theta+2\pi m\ge s+g$
(as $\pi>0$), contradicting $\theta+2\pi k<s+g$. If $k\le m-1$, then
$\theta+2\pi k\le\theta+2\pi m-2\pi\le s$, contradicting $s<\theta+2\pi k$. (The Lean proof makes the same case
split on the sign of $k-m$.)
\end{proof}

\begin{lemma}[shift invariance and relabelling; written step]\label{pc:lem:shift}
Let $s,g,\theta,\alpha,\beta,\delta\in\R$ and $n\in\Z$.
\begin{enumerate}[label=(\alph*)]
\item $\operatorname{inArc}(s,g,\theta+2\pi n)\iff\operatorname{inArc}(s,g,\theta)$.
\item $\operatorname{inArc}(s+2\pi n,g,\theta)\iff\operatorname{inArc}(s,g,\theta)$.
\item $\operatorname{RunMidD}(\alpha,\beta,\delta)\iff\operatorname{RunMidB}(\beta,\delta,\alpha+2\pi)$.
\item $\operatorname{RunMidA}(\alpha,\beta,\delta)\iff\operatorname{RunMidB}(\delta,\alpha+2\pi,\beta+2\pi)$.
\end{enumerate}
\end{lemma}

\begin{proof}
(a) If $s<\theta+2\pi n+2\pi k<s+g$, then $k'=n+k$ witnesses $\operatorname{inArc}(s,g,\theta)$; conversely a
witness $k'$ for $\theta$ gives the witness $k'-n$ for $\theta+2\pi n$.
(b) $s+2\pi n<\theta+2\pi k<s+2\pi n+g$ is equivalent to $s<\theta+2\pi(k-n)<s+g$.

(c) For all $s,g$, $\operatorname{arcHasAntipode}(\beta,\delta,\alpha+2\pi,s,g)$ is the disjunction of
$\operatorname{inArc}(s,g,\beta+\pi)$, $\operatorname{inArc}(s,g,\delta+\pi)$ and
$\operatorname{inArc}(s,g,\alpha+\pi+2\pi)$; by (a) with $n=1$ the last one is equivalent to
$\operatorname{inArc}(s,g,\alpha+\pi)$. Hence
$\operatorname{arcHasAntipode}(\beta,\delta,\alpha+2\pi,s,g)\iff\operatorname{arcHasAntipode}(\alpha,\beta,\delta,s,g)$
for all $s,g$. By definition, $\operatorname{RunMidB}(\beta,\delta,\alpha+2\pi)$ is the conjunction of
$\neg\operatorname{arcHasAntipode}(\beta,\delta,\alpha+2\pi,\beta,\delta-\beta)$ and
$\neg\operatorname{arcHasAntipode}(\beta,\delta,\alpha+2\pi,\delta,(\alpha+2\pi)-\delta)$, which are therefore
equivalent to $\neg H_2$ and $\neg H_3$ (of $(\alpha,\beta,\delta)$). So it is equivalent to
$\operatorname{RunMidD}(\alpha,\beta,\delta)$.

(d) As in (c), $\operatorname{arcHasAntipode}(\delta,\alpha+2\pi,\beta+2\pi,s,g)\iff
\operatorname{arcHasAntipode}(\alpha,\beta,\delta,s,g)$ for all $s,g$ (the antipodes
$\alpha+2\pi+\pi$, $\beta+2\pi+\pi$ are shifted by $2\pi$). $\operatorname{RunMidB}(\delta,\alpha+2\pi,\beta+2\pi)$
is the conjunction of $\neg\operatorname{arcHasAntipode}(\delta,\alpha+2\pi,\beta+2\pi,\delta,(\alpha+2\pi)-\delta)$,
equivalent to $\neg H_3$, and of
$\neg\operatorname{arcHasAntipode}(\delta,\alpha+2\pi,\beta+2\pi,\alpha+2\pi,(\beta+2\pi)-(\alpha+2\pi))$. Since
$(\beta+2\pi)-(\alpha+2\pi)=\beta-\alpha$ and, by (b) with $n=1$,
$\operatorname{inArc}(\alpha+2\pi,\beta-\alpha,\theta)\iff\operatorname{inArc}(\alpha,\beta-\alpha,\theta)$ for
every $\theta$, the second conjunct is equivalent to $\neg H_1$. So the conjunction is equivalent to
$\neg H_3\wedge\neg H_1=\operatorname{RunMidA}(\alpha,\beta,\delta)$.

The Lean proofs of \lean{orient_pos_of_runMidD} and \lean{orient_pos_of_runMidA}
(\dirn{LEANEXT/phase/PhaseClassification.lean}, lines 323--335 and 346--362) prove the implications from left to right
inside the proofs, disjunct by disjunct, keeping the witness $k$ or changing it by $+1$ or $-1$, as (a) and (b)
do above.
\end{proof}

\begin{lemma}[empty arcs; \leanref{PhaseClassification.arc1_empty}{LEANEXT/phase/PhaseClassification.lean:161}, \leanref{PhaseClassification.arc2_empty}{LEANEXT/phase/PhaseClassification.lean:169}, \leanref{PhaseClassification.arc2_empty'}{LEANEXT/phase/PhaseClassification.lean:177}, \leanref{PhaseClassification.arc3_empty}{LEANEXT/phase/PhaseClassification.lean:185}, \leanref{PhaseClassification.arc3_empty'}{LEANEXT/phase/PhaseClassification.lean:193}, \leanref{PhaseClassification.arc1_empty'}{LEANEXT/phase/PhaseClassification.lean:201}]\label{pc:lem:empty}
Let $\alpha,\beta,\delta\in\R$, with $H_1,H_2,H_3$ as in Definition~\ref{pc:def:arc}.
\begin{enumerate}[label=(\alph*)]
\item \textup{(\lean{arc1_empty})} If $\alpha<\beta$, $\beta<\delta$ and $\delta-\alpha\le\pi$, then $\neg H_1$.
\item \textup{(\lean{arc2_empty})} If $\alpha<\beta$, $\beta<\delta$ and $\delta-\alpha\le\pi$, then $\neg H_2$.
\item \textup{(\lean{arc2_empty'})} If $\beta<\delta$, $\delta<\alpha+2\pi$ and $\alpha+2\pi-\beta\le\pi$, then $\neg H_2$.
\item \textup{(\lean{arc3_empty})} If $\beta<\delta$, $\delta<\alpha+2\pi$ and $\alpha+2\pi-\beta\le\pi$, then $\neg H_3$.
\item \textup{(\lean{arc3_empty'})} If $\alpha<\beta$, $\delta<\alpha+2\pi$ and $\beta+2\pi-\delta\le\pi$, then $\neg H_3$.
\item \textup{(\lean{arc1_empty'})} If $\alpha<\beta$, $\delta<\alpha+2\pi$ and $\beta+2\pi-\delta\le\pi$, then $\neg H_1$.
\end{enumerate}
\end{lemma}

\begin{proof}
In each part we apply Lemma~\ref{pc:lem:inarc}(b) to each of the three antipodes $\theta\in\{\alpha+\pi,
\beta+\pi,\delta+\pi\}$ with the integer $m$ indicated, and verify the two inequalities
$s+g\le\theta+2\pi m\le s+2\pi$ for the arc in question; this excludes all three disjuncts of
$\operatorname{arcHasAntipode}$. Throughout we use $\pi>0$. (The Lean proof uses the same integers $m$.)

(a) Arc $H_1$: $s=\alpha$, $s+g=\beta$, $s+2\pi=\alpha+2\pi$. From $\beta<\delta$ and $\delta-\alpha\le\pi$ we get
$\beta<\alpha+\pi$ and $\beta-\alpha<\pi$.
$\theta=\alpha+\pi$, $m=0$: $\beta\le\alpha+\pi\le\alpha+2\pi$.
$\theta=\beta+\pi$, $m=0$: $\beta\le\beta+\pi$, and $\beta+\pi\le\alpha+2\pi$ as $\beta-\alpha<\pi$.
$\theta=\delta+\pi$, $m=0$: $\beta\le\delta+\pi$ as $\beta<\delta$, and $\delta+\pi\le\alpha+2\pi$ as $\delta-\alpha\le\pi$.

(b) Arc $H_2$: $s=\beta$, $s+g=\delta$, $s+2\pi=\beta+2\pi$. From $\alpha<\beta$ and $\delta-\alpha\le\pi$ we get
$\delta-\beta<\pi$.
$\theta=\alpha+\pi$, $m=0$: $\delta\le\alpha+\pi$, and $\alpha+\pi\le\beta+2\pi$ as $\alpha<\beta$.
$\theta=\beta+\pi$, $m=0$: $\delta\le\beta+\pi$ as $\delta-\beta<\pi$, and $\beta+\pi\le\beta+2\pi$.
$\theta=\delta+\pi$, $m=0$: $\delta\le\delta+\pi$, and $\delta+\pi\le\beta+2\pi$ as $\delta-\beta<\pi$.

(c) Arc $H_2$ as in (b). The hypothesis says $\alpha+\pi\le\beta$; with $\delta<\alpha+2\pi$ this gives
$\delta<\beta+\pi$.
$\theta=\alpha+\pi$, $m=1$: $\delta\le\alpha+3\pi$ as $\delta<\alpha+2\pi$, and $\alpha+3\pi\le\beta+2\pi$ as
$\alpha+\pi\le\beta$.
$\theta=\beta+\pi$, $m=0$: $\delta\le\beta+\pi$, and $\beta+\pi\le\beta+2\pi$.
$\theta=\delta+\pi$, $m=0$: $\delta\le\delta+\pi$, and $\delta+\pi\le\beta+2\pi$ as $\delta<\beta+\pi$.

(d) Arc $H_3$: $s=\delta$, $s+g=\alpha+2\pi$, $s+2\pi=\delta+2\pi$. The hypothesis says $\alpha+\pi\le\beta$; with
$\beta<\delta$ this gives $\alpha+\pi<\delta$.
$\theta=\alpha+\pi$, $m=1$: $\alpha+2\pi\le\alpha+3\pi$, and $\alpha+3\pi\le\delta+2\pi$ as $\alpha+\pi<\delta$.
$\theta=\beta+\pi$, $m=0$: $\alpha+2\pi\le\beta+\pi$ as $\alpha+\pi\le\beta$, and $\beta+\pi\le\delta+2\pi$ as
$\beta<\delta$.
$\theta=\delta+\pi$, $m=0$: $\alpha+2\pi\le\delta+\pi$ as $\alpha+\pi<\delta$, and $\delta+\pi\le\delta+2\pi$.

(e) Arc $H_3$ as in (d). The hypothesis says $\beta+\pi\le\delta$; with $\alpha<\beta$ this gives
$\alpha+\pi<\delta$.
$\theta=\alpha+\pi$, $m=1$: $\alpha+2\pi\le\alpha+3\pi$, and $\alpha+3\pi\le\delta+2\pi$ as $\alpha+\pi<\delta$.
$\theta=\beta+\pi$, $m=1$: $\alpha+2\pi\le\beta+3\pi$ as $\alpha<\beta$, and $\beta+3\pi\le\delta+2\pi$ as
$\beta+\pi\le\delta$.
$\theta=\delta+\pi$, $m=0$: $\alpha+2\pi\le\delta+\pi$ as $\alpha+\pi<\delta$, and $\delta+\pi\le\delta+2\pi$.

(f) Arc $H_1$ as in (a). The hypothesis says $\beta+\pi\le\delta$; with $\delta<\alpha+2\pi$ this gives
$\beta<\alpha+\pi$ and $\beta+\pi<\alpha+2\pi$.
$\theta=\alpha+\pi$, $m=0$: $\beta\le\alpha+\pi\le\alpha+2\pi$.
$\theta=\beta+\pi$, $m=0$: $\beta\le\beta+\pi$, and $\beta+\pi\le\alpha+2\pi$.
$\theta=\delta+\pi$, $m=-1$: $\beta\le\delta-\pi$ as $\beta+\pi\le\delta$, and $\delta-\pi\le\alpha+2\pi$ as
$\delta<\alpha+2\pi$.
\end{proof}

\begin{lemma}[a long arc meets an antipode; \leanref{PhaseClassification.arc12_nonempty}{LEANEXT/phase/PhaseClassification.lean:229}]\label{pc:lem:arc12}
Let $\alpha,\beta,\delta\in\R$ with $\alpha<\beta$, $\beta<\delta$, $\delta<\alpha+2\pi$ and $\pi<\delta-\alpha$.
Then $H_1\vee H_2$.
\end{lemma}

\begin{proof}
We give a three-case proof; the Lean proof uses a finer case split (on $\delta-\beta$ against $\pi$ first, then
on $\beta-\alpha$), each case again exhibiting one antipode with an explicit witness $k$.

Case $\beta-\alpha<\pi$: the antipode $\alpha+\pi$ lies in the arc of $H_2$ with $k=0$, since
$\beta<\alpha+\pi$ and $\alpha+\pi<\delta$ (the latter is $\pi<\delta-\alpha$). So $H_2$.

Case $\beta-\alpha>\pi$: $\alpha+\pi$ lies in the arc of $H_1$ with $k=0$, since $\alpha<\alpha+\pi<\beta$. So $H_1$.

Case $\beta-\alpha=\pi$: the antipode $\delta+\pi$ lies in the arc of $H_1$ with $k=-1$: indeed
$\alpha<\delta-\pi$ is $\pi<\delta-\alpha$, and $\delta-\pi<\beta=\alpha+\pi$ is $\delta<\alpha+2\pi$. So $H_1$.
\end{proof}

\subsection{Alternation and runs}

\begin{theorem}[alternation iff all gaps are $<\pi$; \leanref{PhaseClassification.alternating_iff_gaps_lt_pi}{LEANEXT/phase/PhaseClassification.lean:210}]\label{pc:thm:alternating}
Let $\alpha,\beta,\delta\in\R$ with $\alpha<\beta$, $\beta<\delta$ and $\delta<\alpha+2\pi$. Then
\[
  \operatorname{Alternating}(\alpha,\beta,\delta)\iff
  \bigl(\beta-\alpha<\pi\ \wedge\ \delta-\beta<\pi\ \wedge\ \alpha+2\pi-\delta<\pi\bigr).
\]
\end{theorem}

\begin{proof}
($\Rightarrow$) Assume $H_1\wedge H_2\wedge H_3$. If $\beta-\alpha\ge\pi$, then $\alpha+2\pi-\beta\le\pi$, and
Lemma~\ref{pc:lem:empty}(c) (with $\beta<\delta<\alpha+2\pi$) gives $\neg H_2$, a contradiction. If
$\delta-\beta\ge\pi$, then $\beta+2\pi-\delta\le\pi$, and Lemma~\ref{pc:lem:empty}(e) (with $\alpha<\beta$,
$\delta<\alpha+2\pi$) gives $\neg H_3$. If $\alpha+2\pi-\delta\ge\pi$, then $\delta-\alpha\le\pi$, and
Lemma~\ref{pc:lem:empty}(a) gives $\neg H_1$. Hence all three gaps are $<\pi$.

($\Leftarrow$) Assume $g_1,g_2,g_3<\pi$. By Lemma~\ref{pc:lem:inarc}(a):
$H_1$ holds through $\delta+\pi$ with $k=-1$, since $\alpha<\delta-\pi$ (that is $g_3<\pi$) and $\delta-\pi<\beta$
(that is $g_2<\pi$);
$H_2$ holds through $\alpha+\pi$ with $k=0$, since $\beta<\alpha+\pi$ ($g_1<\pi$) and $\alpha+\pi<\delta$
($g_3<\pi$);
$H_3$ holds through $\beta+\pi$ with $k=0$, since $\delta<\beta+\pi$ ($g_2<\pi$) and
$\beta+\pi<\delta+(\alpha+2\pi-\delta)=\alpha+2\pi$ ($g_1<\pi$).
\end{proof}

\begin{remark}[a gap equal to $\pi$]\label{pc:rem:alt-boundary}
No genericity is assumed: a gap may equal $\pi$, and then alternation fails. For example, if $\beta-\alpha=\pi$,
the antipode $\alpha+\pi$ is the endpoint $\beta$ of the open arc $(\beta,\delta)$, and
Lemma~\ref{pc:lem:empty}(c) shows that this arc contains no antipode.
\end{remark}

\begin{proposition}[each arc of an alternating pattern holds exactly the opposite antipode; \leanref{PhaseClassification.alternating_exactly_one}{LEANEXT/phase/PhaseClassification.lean:386}]\label{pc:prop:exactly-one}
Let $\alpha,\beta,\delta\in\R$ with $\alpha<\beta$, $\beta<\delta$, $\delta<\alpha+2\pi$ and
$\operatorname{Alternating}(\alpha,\beta,\delta)$. Then
\begin{enumerate}[label=(\roman*)]
\item $\operatorname{inArc}(\alpha,\beta-\alpha,\delta+\pi)$, $\neg\operatorname{inArc}(\alpha,\beta-\alpha,\alpha+\pi)$ and
  $\neg\operatorname{inArc}(\alpha,\beta-\alpha,\beta+\pi)$;
\item $\operatorname{inArc}(\beta,\delta-\beta,\alpha+\pi)$, $\neg\operatorname{inArc}(\beta,\delta-\beta,\beta+\pi)$ and
  $\neg\operatorname{inArc}(\beta,\delta-\beta,\delta+\pi)$;
\item $\operatorname{inArc}(\delta,\alpha+2\pi-\delta,\beta+\pi)$,
  $\neg\operatorname{inArc}(\delta,\alpha+2\pi-\delta,\alpha+\pi)$ and
  $\neg\operatorname{inArc}(\delta,\alpha+2\pi-\delta,\delta+\pi)$.
\end{enumerate}
\end{proposition}

\begin{proof}
By Theorem~\ref{pc:thm:alternating}, $g_1,g_2,g_3<\pi$. The three memberships were shown in the proof of
Theorem~\ref{pc:thm:alternating}, ($\Leftarrow$). For the non-memberships we use Lemma~\ref{pc:lem:inarc}(b):
\begin{itemize}
\item arc $(\alpha,\beta)$, i.e.\ $s+g=\beta$, $s+2\pi=\alpha+2\pi$: for $\alpha+\pi$ with $m=0$,
  $\beta\le\alpha+\pi$ ($g_1<\pi$) and $\alpha+\pi\le\alpha+2\pi$; for $\beta+\pi$ with $m=0$, $\beta\le\beta+\pi$
  and $\beta+\pi\le\alpha+2\pi$ ($g_1<\pi$);
\item arc $(\beta,\delta)$, i.e.\ $s+g=\delta$, $s+2\pi=\beta+2\pi$: for $\beta+\pi$ with $m=0$,
  $\delta\le\beta+\pi$ ($g_2<\pi$) and $\beta+\pi\le\beta+2\pi$; for $\delta+\pi$ with $m=0$, $\delta\le\delta+\pi$
  and $\delta+\pi\le\beta+2\pi$ ($g_2<\pi$);
\item arc $(\delta,\alpha+2\pi)$, i.e.\ $s+g=\alpha+2\pi$, $s+2\pi=\delta+2\pi$: for $\alpha+\pi$ with $m=1$,
  $\alpha+2\pi\le\alpha+3\pi$ and $\alpha+3\pi\le\delta+2\pi$ ($g_3<\pi$); for $\delta+\pi$ with $m=0$,
  $\alpha+2\pi\le\delta+\pi$ ($g_3<\pi$) and $\delta+\pi\le\delta+2\pi$.\qedhere
\end{itemize}
\end{proof}

\begin{theorem}[a run with middle $\beta$ iff an open half-plane; \leanref{PhaseClassification.runMidB_iff}{LEANEXT/phase/PhaseClassification.lean:251}]\label{pc:thm:runB}
Let $\alpha,\beta,\delta\in\R$ with $\alpha<\beta$, $\beta<\delta$, $\delta<\alpha+2\pi$ and
$\alpha+2\pi-\delta\neq\pi$. Then
\[
  \operatorname{RunMidB}(\alpha,\beta,\delta)\iff\delta-\alpha<\pi .
\]
\end{theorem}

\begin{proof}
($\Leftarrow$) If $\delta-\alpha<\pi$, then $\delta-\alpha\le\pi$, and Lemma~\ref{pc:lem:empty}(a),(b) give
$\neg H_1$ and $\neg H_2$.
($\Rightarrow$) Assume $\neg H_1\wedge\neg H_2$ and, for a contradiction, $\delta-\alpha\ge\pi$. Since
$\alpha+2\pi-\delta\neq\pi$ we have $\delta-\alpha\neq\pi$, hence $\pi<\delta-\alpha$, and
Lemma~\ref{pc:lem:arc12} gives $H_1\vee H_2$, a contradiction.
\end{proof}

\begin{remark}[the gap hypothesis is needed]\label{pc:rem:runB-example}
The direction ($\Rightarrow$) fails without $\alpha+2\pi-\delta\neq\pi$. Take any $\alpha<\beta<\delta$ with
$\delta-\alpha=\pi$, for instance $(\alpha,\beta,\delta)=(0,\pi/2,\pi)$. Then $\delta<\alpha+2\pi$, and
$\operatorname{RunMidB}(\alpha,\beta,\delta)$ holds by Lemma~\ref{pc:lem:empty}(a),(b) (their hypothesis is
$\delta-\alpha\le\pi$), but $\delta-\alpha<\pi$ is false. Geometrically, points at positive radii in the
directions $\alpha$ and $\delta=\alpha+\pi$ lie on opposite rays from $c$, so no open half-plane bounded by a line
through $c$ contains both (compare Proposition~\ref{pc:prop:halfplane}). The direction ($\Leftarrow$) does not use
the gap hypothesis.
\end{remark}

\subsection{Classification under genericity}

\begin{proposition}[gap classification; \leanref{PhaseClassification.gap_classification}{LEANEXT/phase/PhaseClassification.lean:268}]\label{pc:prop:gaps}
Let $\alpha,\beta,\delta\in\R$ with $\alpha<\beta$, $\beta<\delta$, $\delta<\alpha+2\pi$, $\beta-\alpha\neq\pi$,
$\delta-\beta\neq\pi$ and $\alpha+2\pi-\delta\neq\pi$. Then
\[
  \bigl(\beta-\alpha<\pi\wedge\delta-\beta<\pi\wedge\alpha+2\pi-\delta<\pi\bigr)\ \vee\ \delta-\alpha<\pi\ \vee\
  \alpha+2\pi-\beta<\pi\ \vee\ \beta+2\pi-\delta<\pi .
\]
\end{proposition}

\begin{proof}
The gaps $g_1,g_2,g_3$ are positive and sum to $2\pi$, and none equals $\pi$, so each is $<\pi$ or $>\pi$. If all
three are $<\pi$, the first disjunct holds. Two gaps cannot both exceed $\pi$: their sum would exceed $2\pi$,
whereas it equals $2\pi$ minus the (positive) third gap. If exactly one gap exceeds $\pi$, the other two sum to
less than $\pi$: if $g_3>\pi$ then $\delta-\alpha=g_1+g_2=2\pi-g_3<\pi$; if $g_1>\pi$ then
$\alpha+2\pi-\beta=g_2+g_3<\pi$; if $g_2>\pi$ then $\beta+2\pi-\delta=g_3+g_1<\pi$. (The Lean proof runs through
the same $2^3$ sign cases, closing the impossible ones by linear arithmetic.)
\end{proof}

\begin{remark}[boundary cases and exclusivity]\label{pc:rem:gaps}
(1) Genericity cannot be dropped here: for $(\alpha,\beta,\delta)=(0,\pi,3\pi/2)$ all hypotheses hold except
$\beta-\alpha\neq\pi$ (here $g_1=\pi$, $g_2=g_3=\pi/2$), and the four disjuncts read $\pi<\pi$ (first conjunct of the first disjunct), $3\pi/2<\pi$, $\pi<\pi$ and $3\pi/2<\pi$, all
false.

(2) The docstring of \lean{gap_classification} speaks of exactly one case; the Lean statement is only the
disjunction above. The disjuncts are in fact mutually exclusive whenever $\alpha<\beta<\delta<\alpha+2\pi$ (a
written step, not part of the Lean statement): if all gaps are $<\pi$, then each sum of two gaps is
$2\pi$ minus the third, hence $>\pi$; and two of the last three disjuncts together give, after adding,
$2\pi+g_j<2\pi$ for some $j$, which contradicts $g_j>0$.
\end{remark}

\begin{theorem}[phase trichotomy; \leanref{PhaseClassification.phase_trichotomy}{LEANEXT/phase/PhaseClassification.lean:285}]\label{pc:thm:trichotomy}
Let $\alpha,\beta,\delta\in\R$ with $\alpha<\beta$, $\beta<\delta$, $\delta<\alpha+2\pi$, $\beta-\alpha\neq\pi$,
$\delta-\beta\neq\pi$ and $\alpha+2\pi-\delta\neq\pi$. Then
\begin{enumerate}[label=(\roman*)]
\item $\operatorname{Alternating}(\alpha,\beta,\delta)\vee\operatorname{RunMidA}(\alpha,\beta,\delta)\vee
  \operatorname{RunMidB}(\alpha,\beta,\delta)\vee\operatorname{RunMidD}(\alpha,\beta,\delta)$, and
\item $\operatorname{Alternating}(\alpha,\beta,\delta)\Rightarrow\bigl(\neg\operatorname{RunMidA}(\alpha,\beta,\delta)
  \wedge\neg\operatorname{RunMidB}(\alpha,\beta,\delta)\wedge\neg\operatorname{RunMidD}(\alpha,\beta,\delta)\bigr)$.
\end{enumerate}
\end{theorem}

\begin{proof}
(i) Apply Proposition~\ref{pc:prop:gaps}. If all gaps are $<\pi$, Theorem~\ref{pc:thm:alternating} gives
$\operatorname{Alternating}$. If $\delta-\alpha<\pi$, Lemma~\ref{pc:lem:empty}(a),(b) give $\neg H_1\wedge\neg H_2$,
i.e.\ $\operatorname{RunMidB}$. If $\alpha+2\pi-\beta<\pi$, Lemma~\ref{pc:lem:empty}(c),(d) give
$\neg H_2\wedge\neg H_3$, i.e.\ $\operatorname{RunMidD}$. If $\beta+2\pi-\delta<\pi$, Lemma~\ref{pc:lem:empty}(e),(f)
give $\neg H_3\wedge\neg H_1$, i.e.\ $\operatorname{RunMidA}$. (Each strict inequality implies the weak one
required by the lemma.)

(ii) $\operatorname{Alternating}$ is $H_1\wedge H_2\wedge H_3$. Its conjunct $H_3$ contradicts the first conjunct
$\neg H_3$ of $\operatorname{RunMidA}$, $H_1$ contradicts the first conjunct $\neg H_1$ of $\operatorname{RunMidB}$,
and $H_2$ contradicts the first conjunct $\neg H_2$ of $\operatorname{RunMidD}$. This part uses none of the
hypotheses.
\end{proof}

\begin{remark}[what Theorem~\ref{pc:thm:trichotomy} does not state]\label{pc:rem:runs}
The Lean statement is the conjunction of the disjunction (i) and the implication (ii). It does \emph{not} state
that the three runs exclude each other. They do (a written step, not Lean-checked): if
$\alpha<\beta<\delta<\alpha+2\pi$, then $H_1\vee H_2\vee H_3$. Indeed $\alpha<\alpha+\pi<\alpha+2\pi$. If
$\alpha+\pi<\beta$, the antipode $\alpha+\pi$ gives $H_1$ (with $k=0$); if $\beta<\alpha+\pi<\delta$, it gives $H_2$;
if $\delta<\alpha+\pi$, it gives $H_3$, since $\alpha+\pi<\alpha+2\pi$. If $\alpha+\pi=\beta$, then
$\delta\in(\alpha+\pi,\alpha+2\pi)$, so $\delta-\pi\in(\alpha,\beta)$, and $\delta+\pi$ gives $H_1$ with $k=-1$. If
$\alpha+\pi=\delta$, then $\beta\in(\alpha,\alpha+\pi)$, so $\beta+\pi\in(\delta,\alpha+2\pi)$, and $\beta+\pi$ gives
$H_3$ with $k=0$. Any two of $\operatorname{RunMidA}$, $\operatorname{RunMidB}$, $\operatorname{RunMidD}$ together
give $\neg H_1\wedge\neg H_2\wedge\neg H_3$, which is therefore impossible. Hence, under the hypotheses of
Theorem~\ref{pc:thm:trichotomy}, exactly one of the four patterns holds.
\end{remark}

\subsection{The rules in pattern form}

In this subsection $c\in\R^2$, $r_a,r_b,r_d,\alpha,\beta,\delta\in\R$, $\operatorname{polar}$ is
\lean{PhaseClassification.polar}, and
\[
  a=\operatorname{polar}(c,r_a,\alpha),\qquad b=\operatorname{polar}(c,r_b,\beta),\qquad d=\operatorname{polar}(c,r_d,\delta).
\]

\begin{theorem}[rule A from the pattern; \leanref{PhaseClassification.orient_pos_of_alternating}{LEANEXT/phase/PhaseClassification.lean:301}]\label{pc:thm:rule-alt}
If $0<r_a$, $0<r_b$, $0<r_d$, $\alpha<\beta$, $\beta<\delta$, $\delta<\alpha+2\pi$ and
$\operatorname{Alternating}(\alpha,\beta,\delta)$, then $0<\orient(a,b,d)$.
\end{theorem}

\begin{proof}
By Theorem~\ref{pc:thm:alternating}, $\beta-\alpha<\pi$, $\delta-\beta<\pi$ and $\alpha+2\pi-\delta<\pi$; the
ordering gives $0<\beta-\alpha$, $0<\delta-\beta$ and $0<\alpha-\delta+2\pi$. Lemma~\ref{pc:lem:polar-facts}(c)
applies.
\end{proof}

\begin{theorem}[rule B from a run; \leanref{PhaseClassification.orient_pos_of_runMidB}{LEANEXT/phase/PhaseClassification.lean:310}, \leanref{PhaseClassification.orient_pos_of_runMidD}{LEANEXT/phase/PhaseClassification.lean:318}, \leanref{PhaseClassification.orient_pos_of_runMidA}{LEANEXT/phase/PhaseClassification.lean:342}]\label{pc:thm:rule-run}
Let $c$, $r_a,r_b,r_d,\alpha,\beta,\delta$ and $a,b,d$ be as above.
\begin{enumerate}[label=(\Alph*)]
\item[\textup{(B)}] If $0<r_a$, $0<r_d$, $\alpha<\beta$, $\beta<\delta$, $\delta<\alpha+2\pi$,
  $\alpha+2\pi-\delta\neq\pi$, $\operatorname{RunMidB}(\alpha,\beta,\delta)$, $r_a\le r_b$ and $r_d\le r_b$, then
  $0<\orient(a,b,d)$.
\item[\textup{(D)}] If $0<r_a$, $0<r_b$, $\alpha<\beta$, $\beta<\delta$, $\delta<\alpha+2\pi$, $\beta-\alpha\neq\pi$,
  $\operatorname{RunMidD}(\alpha,\beta,\delta)$, $r_b\le r_d$ and $r_a\le r_d$, then $0<\orient(a,b,d)$.
\item[\textup{(A)}] If $0<r_b$, $0<r_d$, $\alpha<\beta$, $\beta<\delta$, $\delta<\alpha+2\pi$, $\delta-\beta\neq\pi$,
  $\operatorname{RunMidA}(\alpha,\beta,\delta)$, $r_b\le r_a$ and $r_d\le r_a$, then $0<\orient(a,b,d)$.
\end{enumerate}
In each case the radius of the middle point of the run is not assumed positive.
\end{theorem}

\begin{proof}
(B) Theorem~\ref{pc:thm:runB} turns $\operatorname{RunMidB}(\alpha,\beta,\delta)$ into $\delta-\alpha<\pi$. With
$0<\beta-\alpha$ and $0<\delta-\beta$ from the ordering, Lemma~\ref{pc:lem:polar-facts}(e) gives the claim.

(D) Put $(\alpha',\beta',\delta')=(\beta,\delta,\alpha+2\pi)$. Then $\alpha'<\beta'$ ($\beta<\delta$),
$\beta'<\delta'$ ($\delta<\alpha+2\pi$) and $\delta'<\alpha'+2\pi$ ($\alpha+2\pi<\beta+2\pi$, from $\alpha<\beta$).
Further $\alpha'+2\pi-\delta'=\beta-\alpha\neq\pi$, and $\operatorname{RunMidB}(\alpha',\beta',\delta')$ holds by
Lemma~\ref{pc:lem:shift}(c). Apply (B) to these directions with the radii $(r_b,r_d,r_a)$ in place of
$(r_a,r_b,r_d)$; its radius hypotheses $0<r_b$, $0<r_a$, $r_b\le r_d$, $r_a\le r_d$ are among ours. It gives
$0<\orient(\operatorname{polar}(c,r_b,\beta),\operatorname{polar}(c,r_d,\delta),\operatorname{polar}(c,r_a,\alpha+2\pi))$.
By Lemma~\ref{pc:lem:periodic}(b) the last point is $a$, so $0<\orient(b,d,a)=\orient(a,b,d)$ by cyclic invariance.

(A) Put $(\alpha',\beta',\delta')=(\delta,\alpha+2\pi,\beta+2\pi)$. Then $\delta<\alpha+2\pi$,
$\alpha+2\pi<\beta+2\pi$ and $\beta+2\pi<\delta+2\pi$, so $\alpha'<\beta'<\delta'<\alpha'+2\pi$; further
$\alpha'+2\pi-\delta'=\delta-\beta\neq\pi$, and $\operatorname{RunMidB}(\alpha',\beta',\delta')$ holds by
Lemma~\ref{pc:lem:shift}(d). Apply (B) with the radii $(r_d,r_a,r_b)$ in place of $(r_a,r_b,r_d)$; its radius
hypotheses $0<r_d$, $0<r_b$, $r_d\le r_a$, $r_b\le r_a$ are among ours. It gives
$0<\orient(\operatorname{polar}(c,r_d,\delta),\operatorname{polar}(c,r_a,\alpha+2\pi),
\operatorname{polar}(c,r_b,\beta+2\pi))=\orient(d,a,b)$ by Lemma~\ref{pc:lem:periodic}(b), and
$\orient(d,a,b)=\orient(a,b,d)$ by cyclic invariance.
\end{proof}

\begin{theorem}[sign from the phase pattern, sorted directions; \leanref{PhaseClassification.orient_pos_of_phase}{LEANEXT/phase/PhaseClassification.lean:372}]\label{pc:thm:phase}
Assume $0<r_a$, $0<r_b$, $0<r_d$, $\alpha<\beta$, $\beta<\delta$, $\delta<\alpha+2\pi$, $\beta-\alpha\neq\pi$,
$\delta-\beta\neq\pi$, $\alpha+2\pi-\delta\neq\pi$, and
\begin{itemize}
\item $\operatorname{RunMidA}(\alpha,\beta,\delta)\Rightarrow(r_b\le r_a\wedge r_d\le r_a)$,
\item $\operatorname{RunMidB}(\alpha,\beta,\delta)\Rightarrow(r_a\le r_b\wedge r_d\le r_b)$,
\item $\operatorname{RunMidD}(\alpha,\beta,\delta)\Rightarrow(r_b\le r_d\wedge r_a\le r_d)$.
\end{itemize}
Then $0<\orient(a,b,d)$.
\end{theorem}

\begin{proof}
By Theorem~\ref{pc:thm:trichotomy}(i) one of four cases holds. If $\operatorname{Alternating}(\alpha,\beta,\delta)$,
apply Theorem~\ref{pc:thm:rule-alt}. If $\operatorname{RunMidA}$, the first implication gives $r_b\le r_a$ and
$r_d\le r_a$, and Theorem~\ref{pc:thm:rule-run}(A) applies (with $0<r_b$, $0<r_d$, $\delta-\beta\neq\pi$). If
$\operatorname{RunMidB}$, the second gives $r_a\le r_b$, $r_d\le r_b$, and (B) applies (with $0<r_a$, $0<r_d$,
$\alpha+2\pi-\delta\neq\pi$). If $\operatorname{RunMidD}$, the third gives $r_b\le r_d$, $r_a\le r_d$, and (D)
applies (with $0<r_a$, $0<r_b$, $\beta-\alpha\neq\pi$).
\end{proof}

The two patterns have a direct geometric meaning.

\begin{proposition}[alternation puts the centre inside the triangle; \leanref{PhaseClassification.centre_inside_of_alternating}{LEANEXT/phase/PhaseClassification.lean:414}]\label{pc:prop:centre}
If $0<r_a$, $0<r_b$, $0<r_d$, $\alpha<\beta$, $\beta<\delta$, $\delta<\alpha+2\pi$ and
$\operatorname{Alternating}(\alpha,\beta,\delta)$, then
\[
  0<\orient(a,b,c),\qquad 0<\orient(b,d,c),\qquad 0<\orient(d,a,c).
\]
\end{proposition}

\begin{proof}
By Theorem~\ref{pc:thm:alternating}, $\beta-\alpha$, $\delta-\beta$ and $\alpha+2\pi-\delta$ lie in $(0,\pi)$. By
Lemma~\ref{pc:lem:polar-facts}(b), applied to the pairs $(a,b)$, $(b,d)$ and $(d,a)$,
\[
  \orient(a,b,c)=r_ar_b\sin(\beta-\alpha),\quad \orient(b,d,c)=r_br_d\sin(\delta-\beta),\quad
  \orient(d,a,c)=r_dr_a\sin(\alpha-\delta),
\]
and $\sin(\alpha-\delta)=\sin(\alpha+2\pi-\delta)$. All three sines and all radii are positive.
\end{proof}

\begin{proposition}[a run with middle $\beta$ lies in an open half-plane through the centre; \leanref{PhaseClassification.halfplane_of_runMidB}{LEANEXT/phase/PhaseClassification.lean:431}]\label{pc:prop:halfplane}
If $0<r_a$, $0<r_b$, $0<r_d$, $\alpha<\beta$, $\beta<\delta$, $\delta<\alpha+2\pi$, $\alpha+2\pi-\delta\neq\pi$ and
$\operatorname{RunMidB}(\alpha,\beta,\delta)$, then there is $\varphi\in\R$ such that for every
$p\in\{a,b,d\}$,
\[
  0<(p_0-c_0)\cos\varphi+(p_1-c_1)\sin\varphi .
\]
\end{proposition}

\begin{proof}
By Theorem~\ref{pc:thm:runB}, $\delta-\alpha<\pi$. Take $\varphi=(\alpha+\delta)/2$. For $p=\operatorname{polar}(c,r,\theta)$,
\[
  (p_0-c_0)\cos\varphi+(p_1-c_1)\sin\varphi=r(\cos\theta\cos\varphi+\sin\theta\sin\varphi)=r\cos(\theta-\varphi).
\]
For $\theta\in\{\alpha,\beta,\delta\}\subseteq[\alpha,\delta]$ we have
$\abs{\theta-\varphi}\le(\delta-\alpha)/2<\pi/2$, so $\cos(\theta-\varphi)>0$; the radius $r\in\{r_a,r_b,r_d\}$ is
positive.
\end{proof}

\subsection{Arbitrary real directions}

\begin{theorem}[sign for arbitrary directions; \leanref{PhaseClassification.orient_sign_of_phase}{LEANEXT/phase/PhaseClassification.lean:496}]\label{pc:thm:sign}
Let $c\in\R^2$ and $r_a,r_b,r_d,\alpha,\beta,\delta\in\R$ with $0<r_a$, $0<r_b$, $0<r_d$, and let
$a=\operatorname{polar}(c,r_a,\alpha)$, $b=\operatorname{polar}(c,r_b,\beta)$, $d=\operatorname{polar}(c,r_d,\delta)$,
where $\operatorname{polar}$ is \lean{PhaseClassification.polar}. Put $\beta'=\operatorname{toIcoMod}(\alpha,\beta)$
and $\delta'=\operatorname{toIcoMod}(\alpha,\delta)$, and assume $\beta'\neq\alpha$ and $\delta'\neq\alpha$. Then
both of the following hold.
\begin{enumerate}[label=(\roman*)]
\item If $\beta'<\delta'$, $\beta'-\alpha\neq\pi$, $\delta'-\beta'\neq\pi$, $\alpha+2\pi-\delta'\neq\pi$,
  $\operatorname{RunMidA}(\alpha,\beta',\delta')\Rightarrow(r_b\le r_a\wedge r_d\le r_a)$,
  $\operatorname{RunMidB}(\alpha,\beta',\delta')\Rightarrow(r_a\le r_b\wedge r_d\le r_b)$ and
  $\operatorname{RunMidD}(\alpha,\beta',\delta')\Rightarrow(r_b\le r_d\wedge r_a\le r_d)$, then $0<\orient(a,b,d)$.
\item If $\delta'<\beta'$, $\delta'-\alpha\neq\pi$, $\beta'-\delta'\neq\pi$, $\alpha+2\pi-\beta'\neq\pi$,
  $\operatorname{RunMidA}(\alpha,\delta',\beta')\Rightarrow(r_d\le r_a\wedge r_b\le r_a)$,
  $\operatorname{RunMidB}(\alpha,\delta',\beta')\Rightarrow(r_a\le r_d\wedge r_b\le r_d)$ and
  $\operatorname{RunMidD}(\alpha,\delta',\beta')\Rightarrow(r_d\le r_b\wedge r_a\le r_b)$, then $\orient(a,b,d)<0$.
\end{enumerate}
\end{theorem}

\begin{proof}
By (T1), $\alpha\le\beta'<\alpha+2\pi$ and $\alpha\le\delta'<\alpha+2\pi$; with $\beta'\neq\alpha$ and
$\delta'\neq\alpha$ we get $\alpha<\beta'$ and $\alpha<\delta'$. By Lemma~\ref{pc:lem:periodic}(c),
$b=\operatorname{polar}(c,r_b,\beta')$ and $d=\operatorname{polar}(c,r_d,\delta')$.

(i) The directions $(\alpha,\beta',\delta')$ satisfy $\alpha<\beta'<\delta'<\alpha+2\pi$, and the remaining
hypotheses of (i) are exactly the gap conditions and radius implications of Theorem~\ref{pc:thm:phase} for these
directions and the radii $(r_a,r_b,r_d)$. Hence $0<\orient(a,b,d)$.

(ii) Apply Theorem~\ref{pc:thm:phase} to the directions $(\alpha,\delta',\beta')$, which satisfy
$\alpha<\delta'<\beta'<\alpha+2\pi$, and the radii $(r_a,r_d,r_b)$ in place of $(r_a,r_b,r_d)$. Its gap conditions
are $\delta'-\alpha\neq\pi$, $\beta'-\delta'\neq\pi$, $\alpha+2\pi-\beta'\neq\pi$, and its radius implications,
with $r_b$ and $r_d$ interchanged, are exactly the three implications of (ii). Hence
$0<\orient(a,\operatorname{polar}(c,r_d,\delta'),\operatorname{polar}(c,r_b,\beta'))=\orient(a,d,b)=-\orient(a,b,d)$,
so $\orient(a,b,d)<0$.
\end{proof}

\begin{remark}[scope of Theorem~\ref{pc:thm:sign}]\label{pc:rem:sign-scope}
The case $\beta'=\delta'$ (the points $b$ and $d$ on one ray from $c$) is not covered, and $\beta'=\alpha$ or
$\delta'=\alpha$ is excluded by hypothesis. In Lean this theorem is not combined with the layer theorems of
Subsection~\ref{pc:sub:layers} (\dirn{LEANEXT/capstone/REPORT.md}, open issues).
\end{remark}

\subsection{The discrete dichotomy of cyclic words}

Let $\mathrm{Fin}\,6=\{0,1,\dots,5\}$ with addition modulo $6$, and $\mathrm{Bool}=\{\mathrm{true},\mathrm{false}\}$
with negation $\neg$. A Boolean value $x$ used as a statement means $x=\mathrm{true}$.

\begin{theorem}[discrete dichotomy, computer-assisted in Lean; written proof by hand; \leanref{PhaseClassification.word_dichotomy}{LEANEXT/phase/PhaseClassification.lean:539}]\label{pc:thm:word}
For every $s\colon\mathrm{Fin}\,6\to\mathrm{Bool}$ such that $s(i+3)=\neg s(i)$ for all $i\in\mathrm{Fin}\,6$:
\begin{enumerate}[label=(\roman*)]
\item $\bigl(\forall i:\ s(i+1)=\neg s(i)\bigr)\ \vee\ \bigl(\exists i:\ s(i)\wedge s(i+1)\wedge s(i+2)\bigr)$, and
\item $\neg\Bigl(\bigl(\forall i:\ s(i+1)=\neg s(i)\bigr)\wedge\bigl(\exists i:\ s(i)\wedge s(i+1)\wedge s(i+2)\bigr)\Bigr)$.
\end{enumerate}
\end{theorem}

\paragraph{How Lean checks it.} The statement quantifies over the finite type of functions
$\mathrm{Fin}\,6\to\mathrm{Bool}$, which has $2^6=64$ elements, and inside over $\mathrm{Fin}\,6$; all atoms are
equalities of Booleans. Lean therefore has a decision procedure for it. For each of the $64$ functions $s$, the
procedure decides the implication from the premise (the six equalities $s(i+3)=\neg s(i)$) to the conclusion
(i)~$\wedge$~(ii), whose atoms are the six equalities $s(i+1)=\neg s(i)$ of the alternation condition and the six
triples $s(i),s(i+1),s(i+2)$ of the run condition. Lean's decision procedure for an implication evaluates the
premise first and the conclusion only when the premise holds, so the conclusion is evaluated only for the $8$
functions listed in the proof below. The tactic \lean{decide} evaluates the procedure during elaboration, and the
Lean kernel evaluates it again, to \lean{true}, when it type-checks the resulting proof term
(\lean{of_decide_eq_true} applied to \lean{Eq.refl}~\lean{true}). No
\lean{native_decide} is used, so no \lean{Lean.ofReduceBool} axiom enters; \dirn{LEANEXT/phase/logs/axioms.log}
records the three standard axioms. Mutant M7 of the mutation script \dirn{LEANEXT/phase/mutants.py} replaces the
premise by \lean{True}; the mutation log \dirn{LEANEXT/phase/logs/mutants.log}, which calls this mutant ``drop
antipodal premise in \lean{word_dichotomy}'', records that it fails to compile with the error ``Tactic
\lean{decide} proved that the proposition'' (the script logs only the first line of the first error message;
Lean's message continues with the proposition and ``is false''). The mutated statement is indeed false: for
example, the word $PPAPPA$ (with $P=\mathrm{true}$, $A=\mathrm{false}$; it violates the premise, as $s(3)=s(0)$) is
neither alternating nor contains a cyclic $PPP$.

\begin{proof}[Written proof (by hand)]
Let $s$ satisfy the premise and put $x=s(0)$, $y=s(1)$, $z=s(2)$. The premise gives $s(3)=\neg x$, $s(4)=\neg y$,
$s(5)=\neg z$, so $s=(x,y,z,\neg x,\neg y,\neg z)$ and there are exactly eight such words. Conversely each of
these eight words satisfies the premise, since $s(i+3)=\neg s(i)$ for $i=0,1,2$ gives it for $i=3,4,5$ as well. In
the notation $P=\mathrm{true}$, $A=\mathrm{false}$, and listing positions $0,\dots,5$, they are:
\begin{center}
\begin{tabular}{lll}
\toprule
$(x,y,z)$ & word & property \\
\midrule
$(P,P,P)$ & $PPPAAA$ & run $PPP$ at $i=0$ \\
$(P,P,A)$ & $PPAAAP$ & run at $i=5$ (positions $5,0,1$) \\
$(P,A,P)$ & $PAPAPA$ & alternating \\
$(P,A,A)$ & $PAAAPP$ & run at $i=4$ (positions $4,5,0$) \\
$(A,P,P)$ & $APPPAA$ & run at $i=1$ \\
$(A,P,A)$ & $APAPAP$ & alternating \\
$(A,A,P)$ & $AAPPPA$ & run at $i=2$ \\
$(A,A,A)$ & $AAAPPP$ & run at $i=3$ \\
\bottomrule
\end{tabular}
\end{center}
This proves (i). For (ii): if $s(i+1)=\neg s(i)$ for all $i$, then $s(i)$ and $s(i+1)$ are never both true, so
there is no cyclic run. (A shorter argument for (i): if $s$ is not alternating, then $s(i+1)=s(i)$ for some $i$, and
then also $s(i+4)=\neg s(i+1)=\neg s(i)=s(i+3)$; one of the equal pairs $\{i,i+1\}$, $\{i+3,i+4\}$ consists of two
$P$'s, say at $j,j+1$. If $s(j+2)=P$ there is a run at $j$; otherwise $s(j-1)=s(j+5)=\neg s(j+2)=P$ and there is a
run at $j-1$.)
\end{proof}

\begin{remark}[scope of Theorem~\ref{pc:thm:word}]\label{pc:rem:word-scope}
This is a statement about Boolean words only. The layered encoders read the six directions $\alpha,\beta,\delta,
\alpha+\pi,\beta+\pi,\delta+\pi$, sorted counterclockwise, as such a word ($P$ for a direction, $A$ for an
antipode). That this sorted word satisfies $s(i+3)=\neg s(i)$, and that a cyclic $PPP$ in it is
$\operatorname{RunMidA}$, $\operatorname{RunMidB}$ or $\operatorname{RunMidD}$ of the middle point, is not proved in
\dirn{LEANEXT} (\dirn{LEANEXT/phase/REPORT.md}, ``Not in Lean'', item 2). For encoder~B, the Lean version of the
encoder in the ENCBLEAN package has a Lean-checked counterpart of the second link that does not use
Theorem~\ref{pc:thm:word}: \lean{EncB.wordSpec} (\dirn{ENCBLEAN/EncB/SoundWord.lean:222}) derives the rule-B
pattern \lean{EncB.BlockCCW} from three cyclically consecutive points of the encoder's sorted word, and, in
\dirn{ENCBLEAN/geo/ES8Geo/ArcBridge.lean}, \lean{EncB.Geo.Arc.arcFrame} (line 393) identifies, for a polar labelling
of the layers, the arc tests of this pattern with the tests $\operatorname{arcHasAntipode}$ for true directions
$\alpha<\beta<\delta<\alpha+2\pi$, from which \lean{EncB.Geo.labelling_ruleB} (line 473) obtains
$\operatorname{RunMidB}(\alpha,\beta,\delta)$ and applies Theorem~\ref{pc:thm:rule-run}(B). These three
declarations depend only on \lean{propext}, \lean{Classical.choice} and \lean{Quot.sound}: for \lean{wordSpec} by
the \texttt{\#guard\_msgs} pin that follows it (\dirn{ENCBLEAN/EncB/SoundWord.lean}, lines 228--229; the module is
built in \dirn{ENCBLEAN/logs/m3_build.log}, line 9, and \dirn{ENCBLEAN/logs/axioms.log}, line 446, gives the same
list), and for \lean{arcFrame} and \lean{labelling_ruleB} by their \texttt{\#print axioms} output in
\dirn{ENCBLEAN/logs/m8_fix_build.log}, lines 35 and 29, and by the pin of \lean{EncB.Geo.modelOfLabelling} (which
uses \lean{labelling_ruleB}) described under Rigor and axioms. The case split for real directions that is
Lean-checked here is Theorem~\ref{pc:thm:trichotomy}.
\end{remark}

\subsection{Rule B on the layers of a symmetric set}\label{pc:sub:layers}

In this subsection $S$ is a finite set of points, $c\in\R^2$, $m,K\in\N$, and $\rho=\operatorname{rotAbout}(c,2\pi/m)$
is the counterclockwise rotation about $c$ by $2\pi/m$ (Definition~\ref{fd:def:rot}; $m$ is cast to $\R$). We say that
\emph{the layer hypotheses \textup{(L)} hold} if
\[
  0<m,\qquad K\le 2m,\qquad c\notin S,\qquad \operatorname{Invariant}(\rho,S),\qquad \neg\operatorname{HasConvexNGon}(K,S),
\]
where $\operatorname{Invariant}(\rho,S)$ means $\rho(p)\in S$ for all $p\in S$ (Definition~\ref{fd:def:inv}) and
$\operatorname{HasConvexNGon}(K,S)$ means that some $K$-element subset of $S$ is in convex position
(Definition~\ref{fd:def:ngon}). In this subsection $\operatorname{polar}$ is \lean{LayerLemma.polar}; by
Lemma~\ref{pc:lem:helpers}(a) it is the same function as \lean{PhaseClassification.polar}, so the results of the
previous subsections apply to it.

We use two facts from the layer sections. First, $\peel(S,n)\subseteq S$ and $\hv(T)\subseteq T$ for all $n$ and
all finite $T$ (Lean \lean{peel_subset}, \lean{hullVerts_subset}): $\hv(T)$ is defined as a subset of $T$
(Definition~\ref{ll:def:hv}), and $\peel(S,n+1)=\peel(S,n)\setminus\hv(\peel(S,n))\subseteq\peel(S,n)$
(Definition~\ref{ll:def:peel}), so induction on $n$ gives the first inclusion. Second, the radius order of the
layers (Theorem~\ref{lg:thm:radius}; Lean \lean{ErdosSzekeres.LayerLemma.dist2_le_of_later},
\dirn{LAYERLEAN/LayerLemma.lean:1012}, repeated at \dirn{LEANEXT/linked/Linked.lean:1012} and
\dirn{LEANEXT/capstone/Capstone2.lean:1020}): \emph{if \textup{(L)} holds, $n\le n'$, $x\in\hv(\peel(S,n))$ and
$y\in\peel(S,n')$, then $\operatorname{dist2}(c,y)\le\operatorname{dist2}(c,x)$.} We also use the elementary fact
that $0<r$, $0<r'$ and $r^2\le r'^2$ imply $r\le r'$ (otherwise $r^2-r'^2=(r-r')(r+r')>0$).

\begin{theorem}[rule B on the convex layers; \leanref{ErdosSzekeres.LayerLemma.orient_pos_of_layers}{LEANEXT/linked/Linked.lean:1398} for (a), and the six declarations of \lean{ErdosSzekeres.Capstone} cited in (b)--(e)]\label{pc:thm:layers}
Assume \textup{(L)}. Let $n_a,n_b,n_d\in\N$ and $a,b,d\in\R^2$. For real $r_a,r_b,r_d,\alpha,\beta,\delta$ let
\textup{(P)} denote the conditions
\[
  0<r_a,\quad 0<r_b,\quad 0<r_d,\quad a=\operatorname{polar}(c,r_a,\alpha),\quad b=\operatorname{polar}(c,r_b,\beta),
  \quad d=\operatorname{polar}(c,r_d,\delta),
\]
and for real $\alpha,\beta,\delta$ let \textup{(P$\surd$)} denote the conditions
\[
  a=\operatorname{polar}\bigl(c,\sqrt{\operatorname{dist2}(c,a)},\alpha\bigr),\quad
  b=\operatorname{polar}\bigl(c,\sqrt{\operatorname{dist2}(c,b)},\beta\bigr),\quad
  d=\operatorname{polar}\bigl(c,\sqrt{\operatorname{dist2}(c,d)},\delta\bigr).
\]
\begin{enumerate}[label=(\alph*)]
\item \textup{(\lean{orient_pos_of_layers})} If $a\in\peel(S,n_a)$, $b\in\hv(\peel(S,n_b))$, $d\in\peel(S,n_d)$,
  $n_b\le n_a$, $n_b\le n_d$, and $r_a,r_b,r_d,\alpha,\beta,\delta\in\R$ satisfy \textup{(P)}, $0<\beta-\alpha$,
  $0<\delta-\beta$ and $\delta-\alpha<\pi$, then $0<\orient(a,b,d)$.
\item \textup{(\leanref{ErdosSzekeres.Capstone.orient_pos_of_layers_runMidB}{LEANEXT/capstone/Capstone.lean:1977}.)}
  If $a\in\peel(S,n_a)$, $b\in\hv(\peel(S,n_b))$,
  $d\in\peel(S,n_d)$, $n_b\le n_a$, $n_b\le n_d$, and $r_a,r_b,r_d,\alpha,\beta,\delta\in\R$ satisfy \textup{(P)},
  $\alpha<\beta$, $\beta<\delta$, $\delta<\alpha+2\pi$, $\alpha+2\pi-\delta\neq\pi$ and
  $\operatorname{RunMidB}(\alpha,\beta,\delta)$, then $0<\orient(a,b,d)$.
\item \textup{(\leanref{ErdosSzekeres.Capstone.orient_pos_of_layers_runMidD}{LEANEXT/capstone/Capstone2.lean:2011}.)}
  If $a\in\peel(S,n_a)$, $b\in\peel(S,n_b)$,
  $d\in\hv(\peel(S,n_d))$, $n_d\le n_a$, $n_d\le n_b$, and $r_a,r_b,r_d,\alpha,\beta,\delta\in\R$ satisfy
  \textup{(P)}, $\alpha<\beta$, $\beta<\delta$, $\delta<\alpha+2\pi$, $\beta-\alpha\neq\pi$ and
  $\operatorname{RunMidD}(\alpha,\beta,\delta)$, then $0<\orient(a,b,d)$.
\item \textup{(\leanref{ErdosSzekeres.Capstone.orient_pos_of_layers_runMidA}{LEANEXT/capstone/Capstone2.lean:2031}.)}
  If $a\in\hv(\peel(S,n_a))$, $b\in\peel(S,n_b)$,
  $d\in\peel(S,n_d)$, $n_a\le n_b$, $n_a\le n_d$, and $r_a,r_b,r_d,\alpha,\beta,\delta\in\R$ satisfy \textup{(P)},
  $\alpha<\beta$, $\beta<\delta$, $\delta<\alpha+2\pi$, $\delta-\beta\neq\pi$ and
  $\operatorname{RunMidA}(\alpha,\beta,\delta)$, then $0<\orient(a,b,d)$.
\item \textup{(The \lean{_intrinsic} versions.)} Let $\alpha,\beta,\delta\in\R$. Then $0<\orient(a,b,d)$ in each
  of the following three cases:
  \begin{enumerate}[label=(e\arabic*)]
  \item \textup{(\leanref{ErdosSzekeres.Capstone.orient_pos_of_layers_runMidB_intrinsic}{LEANEXT/capstone/Capstone2.lean:2082}.)}
    $a\in\peel(S,n_a)$, $b\in\hv(\peel(S,n_b))$,
    $d\in\peel(S,n_d)$, $n_b\le n_a$, $n_b\le n_d$, \textup{(P$\surd$)}, $\alpha<\beta$, $\beta<\delta$,
    $\delta<\alpha+2\pi$, $\alpha+2\pi-\delta\neq\pi$ and $\operatorname{RunMidB}(\alpha,\beta,\delta)$;
  \item \textup{(\leanref{ErdosSzekeres.Capstone.orient_pos_of_layers_runMidD_intrinsic}{LEANEXT/capstone/Capstone2.lean:2098}.)}
    $a\in\peel(S,n_a)$, $b\in\peel(S,n_b)$,
    $d\in\hv(\peel(S,n_d))$, $n_d\le n_a$, $n_d\le n_b$, \textup{(P$\surd$)}, $\alpha<\beta$, $\beta<\delta$,
    $\delta<\alpha+2\pi$, $\beta-\alpha\neq\pi$ and $\operatorname{RunMidD}(\alpha,\beta,\delta)$;
  \item \textup{(\leanref{ErdosSzekeres.Capstone.orient_pos_of_layers_runMidA_intrinsic}{LEANEXT/capstone/Capstone2.lean:2115}.)}
    $a\in\hv(\peel(S,n_a))$, $b\in\peel(S,n_b)$,
    $d\in\peel(S,n_d)$, $n_a\le n_b$, $n_a\le n_d$, \textup{(P$\surd$)}, $\alpha<\beta$, $\beta<\delta$,
    $\delta<\alpha+2\pi$, $\delta-\beta\neq\pi$ and $\operatorname{RunMidA}(\alpha,\beta,\delta)$.
  \end{enumerate}
\end{enumerate}
General position of $S$ is not assumed.
\end{theorem}

\begin{proof}
(a) Apply the radius order with $(n,n',x,y)=(n_b,n_a,b,a)$ and with $(n_b,n_d,b,d)$: this gives
$\operatorname{dist2}(c,a)\le\operatorname{dist2}(c,b)$ and $\operatorname{dist2}(c,d)\le\operatorname{dist2}(c,b)$. By
(P) and Lemma~\ref{pc:lem:helpers}(c) these read $r_a^2\le r_b^2$ and $r_d^2\le r_b^2$, hence $r_a\le r_b$ and
$r_d\le r_b$ as all radii are positive. Rule~B, Lemma~\ref{pc:lem:polar-facts}(e) (in Lean the \lean{LayerLemma}
copy, Lemma~\ref{lg:lem:ruleB}), with $0<r_a$, $0<r_d$, $0<\beta-\alpha$, $0<\delta-\beta$ and $\delta-\alpha<\pi$
gives $0<\orient(a,b,d)$.

(b) By Theorem~\ref{pc:thm:runB} (with $\alpha<\beta<\delta<\alpha+2\pi$ and $\alpha+2\pi-\delta\neq\pi$),
$\operatorname{RunMidB}(\alpha,\beta,\delta)$ gives $\delta-\alpha<\pi$. Together with $0<\beta-\alpha$ and
$0<\delta-\beta$, part (a) applies.

(c) The radius order with $(n_d,n_a,d,a)$ and $(n_d,n_b,d,b)$ gives $\operatorname{dist2}(c,a)\le\operatorname{dist2}(c,d)$
and $\operatorname{dist2}(c,b)\le\operatorname{dist2}(c,d)$, so $r_a^2\le r_d^2$ and $r_b^2\le r_d^2$, hence
$r_a\le r_d$ and $r_b\le r_d$. Theorem~\ref{pc:thm:rule-run}(D), with $0<r_a$, $0<r_b$, the ordering,
$\beta-\alpha\neq\pi$, $\operatorname{RunMidD}(\alpha,\beta,\delta)$, $r_b\le r_d$ and $r_a\le r_d$, gives
$0<\orient(a,b,d)$.

(d) The radius order with $(n_a,n_b,a,b)$ and $(n_a,n_d,a,d)$ gives $r_b^2\le r_a^2$ and $r_d^2\le r_a^2$, hence
$r_b\le r_a$ and $r_d\le r_a$. Theorem~\ref{pc:thm:rule-run}(A), with $0<r_b$, $0<r_d$, the ordering,
$\delta-\beta\neq\pi$, $\operatorname{RunMidA}(\alpha,\beta,\delta)$, $r_b\le r_a$ and $r_d\le r_a$, gives
$0<\orient(a,b,d)$.

(e) In each case $a,b,d\in S$: a point of $\peel(S,n)$ lies in $S$, and a point of $\hv(\peel(S,n))$ lies in
$\peel(S,n)\subseteq S$. As $c\notin S$, Lemma~\ref{pc:lem:helpers}(f) gives $a,b,d\neq c$, and
Lemma~\ref{pc:lem:helpers}(e) gives $0<\sqrt{\operatorname{dist2}(c,a)}$, $0<\sqrt{\operatorname{dist2}(c,b)}$ and
$0<\sqrt{\operatorname{dist2}(c,d)}$. So (P) holds with $r_a=\sqrt{\operatorname{dist2}(c,a)}$,
$r_b=\sqrt{\operatorname{dist2}(c,b)}$, $r_d=\sqrt{\operatorname{dist2}(c,d)}$, and (e1), (e2), (e3) follow from (b),
(c), (d) respectively.
\end{proof}

\begin{corollary}[rule A restated; \leanref{ErdosSzekeres.Capstone.orient_pos_of_alternating_polar}{LEANEXT/capstone/Capstone.lean:1992}]\label{pc:cor:alt-polar}
Let $c\in\R^2$ and $r_a,r_b,r_d,\alpha,\beta,\delta\in\R$ with $0<r_a$, $0<r_b$, $0<r_d$, $\alpha<\beta$,
$\beta<\delta$, $\delta<\alpha+2\pi$ and $\operatorname{Alternating}(\alpha,\beta,\delta)$. Then, with
$\operatorname{polar}$ the constant \lean{LayerLemma.polar},
$0<\orient(\operatorname{polar}(c,r_a,\alpha),\operatorname{polar}(c,r_b,\beta),\operatorname{polar}(c,r_d,\delta))$.
No layer hypotheses are involved.
\end{corollary}

\begin{proof}
By Lemma~\ref{pc:lem:helpers}(a) this is Theorem~\ref{pc:thm:rule-alt}. (The Lean proof repeats the proof of
Theorem~\ref{pc:thm:rule-alt} with the \lean{LayerLemma} copy of rule~A, Lemma~\ref{lg:lem:ruleA}.)
\end{proof}

\begin{remark}[what the layer theorems assume]\label{pc:rem:layers-scope}
In Theorem~\ref{pc:thm:layers} the polar forms (with radius variables in (a)--(d), with the intrinsic radii
$\sqrt{\operatorname{dist2}(c,\cdot)}$ in (e)) and the conditions on the given directions are hypotheses: in (a)
the conditions $0<\beta-\alpha$, $0<\delta-\beta$ and $\delta-\alpha<\pi$; in (b)--(e) the strict order
$\alpha<\beta<\delta<\alpha+2\pi$, the gap condition $\neq\pi$ and the pattern. They are not derived from the
symmetry or from general position. Directions with (P$\surd$) exist for every point other than $c$
(Lemma~\ref{pc:lem:helpers}(d)), but choosing them in sorted order and excluding gaps equal to $\pi$ is left to the
user of these theorems; the reduction of arbitrary directions in Theorem~\ref{pc:thm:sign} is not combined with
the layer hypotheses in Lean. The middle point (the point $b$ of the middle direction in (a), the middle point of
the run in (b)--(e)) is required to be a hull vertex of a layer whose index is at most the indices of the layers
of the other two points, and the conclusion is about the fixed order $(a,b,d)$.
\end{remark}

\paragraph{Coverage.} All declarations of \dirn{LEANEXT/phase/PhaseClassification.lean} are covered:
\lean{polar} (Definition~\ref{pc:def:polar}); \lean{orient_polar}, \lean{orient_polar_centre},
\lean{orient_polar_centre_inside}, \lean{sin_add_lt}, \lean{orient_polar_halfplane}
(Lemma~\ref{pc:lem:polar-facts}; except \lean{orient_polar_centre}, these are the copies of lines 931--996 of
\dirn{LAYERLEAN/LayerLemma.lean}, whose originals also have written proofs in the section on layer geometry);
\lean{polar_add_int_two_pi},
\lean{polar_add_two_pi}, \lean{polar_toIcoMod} (Lemma~\ref{pc:lem:periodic}); \lean{exists_polar}
(Lemma~\ref{pc:lem:exists-polar}); \lean{inArc}, \lean{arcHasAntipode}, \lean{Alternating}, \lean{RunMidB},
\lean{RunMidD}, \lean{RunMidA} (Definition~\ref{pc:def:arc}); \lean{inArc_of}, \lean{not_inArc}
(Lemma~\ref{pc:lem:inarc}); \lean{arc1_empty}, \lean{arc2_empty}, \lean{arc2_empty'}, \lean{arc3_empty},
\lean{arc3_empty'}, \lean{arc1_empty'} (Lemma~\ref{pc:lem:empty}); \lean{arc12_nonempty}
(Lemma~\ref{pc:lem:arc12}); \lean{alternating_iff_gaps_lt_pi} (Theorem~\ref{pc:thm:alternating});
\lean{alternating_exactly_one} (Proposition~\ref{pc:prop:exactly-one}); \lean{runMidB_iff}
(Theorem~\ref{pc:thm:runB}); \lean{gap_classification} (Proposition~\ref{pc:prop:gaps}); \lean{phase_trichotomy}
(Theorem~\ref{pc:thm:trichotomy}); \lean{orient_pos_of_alternating} (Theorem~\ref{pc:thm:rule-alt});
\lean{orient_pos_of_runMidB}, \lean{orient_pos_of_runMidD}, \lean{orient_pos_of_runMidA}
(Theorem~\ref{pc:thm:rule-run}); \lean{orient_pos_of_phase} (Theorem~\ref{pc:thm:phase});
\lean{centre_inside_of_alternating} (Proposition~\ref{pc:prop:centre}); \lean{halfplane_of_runMidB}
(Proposition~\ref{pc:prop:halfplane}); \lean{orient_sign_of_phase} (Theorem~\ref{pc:thm:sign});
\lean{word_dichotomy} (Theorem~\ref{pc:thm:word}, computer-assisted in Lean; written proof by hand). The new declarations of
\dirn{LEANEXT/linked/Linked.lean} are covered: \lean{dist2_polar} (Lemma~\ref{pc:lem:helpers}(c)) and
\lean{orient_pos_of_layers} (Theorem~\ref{pc:thm:layers}(a)); so are those of \dirn{LEANEXT/capstone/Capstone.lean}:
\lean{orient_pos_of_layers_runMidB} (Theorem~\ref{pc:thm:layers}(b)), \lean{orient_pos_of_alternating_polar}
(Corollary~\ref{pc:cor:alt-polar}) and \lean{polar_eq} (Lemma~\ref{pc:lem:helpers}(a)); and those of
\dirn{LEANEXT/capstone/Capstone2.lean}: \lean{orient_pos_of_layers_runMidD}, \lean{orient_pos_of_layers_runMidA}
(Theorem~\ref{pc:thm:layers}(c),(d)), \lean{polar_zero}, \lean{exists_angle}, \lean{sqrt_dist2_pos},
\lean{ne_of_mem} (Lemma~\ref{pc:lem:helpers}(b),(d),(e),(f)) and the three \lean{_intrinsic} theorems
(Theorem~\ref{pc:thm:layers}(e)). Deliberately not covered separately: lines 1--1384 of \dirn{Linked.lean} and the
corresponding lines of \dirn{Capstone.lean} and \dirn{Capstone2.lean} (copies of \dirn{LAYERLEAN/LayerLemma.lean}
and of \dirn{LEANEXT/merged/Merged.lean}, treated in the sections on the layer lemma, layer geometry and orbit
literals), and lines 1426--1965 of \dirn{Capstone.lean} and \dirn{Capstone2.lean} (verbatim copies of
\dirn{PhaseClassification.lean}, whose statements and proofs are those covered above, at line offset $+1420$).
The helper Lemma~\ref{pc:lem:shift} is a written step that makes explicit two \lean{have} steps of the Lean proofs
of \lean{orient_pos_of_runMidD} and \lean{orient_pos_of_runMidA}; Remarks~\ref{pc:rem:alt-boundary}, \ref{pc:rem:runB-example}, \ref{pc:rem:gaps} and~\ref{pc:rem:runs} contain
further written steps (with proofs) that are not Lean statements.
\par\endgroup
